\documentclass[11pt]{article}
\usepackage[margin=1.05in]{geometry}
\usepackage{amsmath,amssymb,amsthm}
\usepackage{booktabs}
\usepackage{microtype}
\usepackage{xcolor}
\usepackage[colorlinks=true,linkcolor=blue!55!black,citecolor=blue!55!black,urlcolor=blue!55!black]{hyperref}
\usepackage{mathtools}
\usepackage{enumitem}
\setlist{topsep=3pt,itemsep=2pt}

\newcommand{\eps}{\varepsilon}
\newcommand{\dd}{\mathrm{d}}
\newcommand{\Pexp}{\mathbb{P}\exp}
\newcommand{\Q}{\mathbb{Q}}
\newcommand{\Z}{\mathbb{Z}}
\newcommand{\Fp}{\mathbb{F}_p}
\newtheorem{definition}{Definition}[section]
\newtheorem{example}[definition]{Example}
\newtheorem{remark}[definition]{Remark}

\title{Differential Equations for Master Integrals\\[0.4ex]
\large Lecture notes on the analytic evaluation of the double-real
phase-space integrals in $pp \to h + X$ at NNLO}
\date{August 24, 2026}

\begin{document}
\maketitle

\begin{abstract}
\noindent
Integration-by-parts reduction expresses all phase-space integrals of
the double-real channel through a finite basis of master integrals.
These notes develop, from first principles, the method by which the
masters are evaluated: the derivation of first-order differential
equations in the kinematic invariants; the organization of the
equations into triangular systems; the identification of integral
families by changes of loop momenta and, separately, of the diagonal
blocks of the resulting systems by conjugation, which together reduce
the many systems to a short list of genuinely distinct ones; the
canonical form of the equations, in which the Laurent expansion in the
dimensional regulator integrates order by order into iterated
integrals; the treatment of square roots of kinematic polynomials by
rational changes of variables and, where no such change of variables
exists, by the Galois theory of the roots; the determination of
boundary conditions from the analytic and geometric properties of the
phase space; and the recovery of the distributional behaviour of the
integrals at the edges of phase space.
\end{abstract}

\setcounter{tocdepth}{2}
\tableofcontents

%======================================================================
\section{Introduction}
\label{sec:setting}

\subsection{Phase-space integrals as cut loop integrals}

The physical quantity behind these notes is the NNLO hard function
for single-inclusive hadron production, $pp \to h + X$, in collinear
factorization.  In its double-real channel two unresolved massless
partons are radiated and integrated over, so the basic objects are
phase-space integrals of squared tree amplitudes.  It is convenient
to convert them into loop integrals at the outset.  By the
Sokhotski--Plemelj identity,
\begin{equation}
  \delta^{(+)}(p^2) \;=\;
  \frac{1}{2\pi i}\left[\frac{1}{p^2 - i0} - \frac{1}{p^2 + i0}\right]
  \quad\text{on positive-energy configurations},
\end{equation}
each on-shell delta function is the discontinuity of a propagator;
the positive-energy condition $\theta(p^0)$ is supplied by the cut
itself rather than by the discontinuity.  Every phase-space integral
thereby becomes a two-loop integral in which certain lines are
\emph{cut}.  Cut integrals obey the same integration-by-parts
identities as ordinary loop integrals, because the delta functions are
ordinary distributions: differentiation under the integral sign and
the vanishing of total derivatives are unimpaired.  This rewriting is
the device known as reverse unitarity~\cite{am}.  Two rules govern a cut line.
Raising its power to $a_i \ge 1$ produces
\begin{equation}
  \frac{1}{D_i^{\,a_i}} \;\longleftrightarrow\;
  \frac{(-1)^{a_i-1}}{(a_i-1)!}\,
  \delta^{(a_i-1)}\big(D_i\big) ,
  \label{eq:cutpower}
\end{equation}
again an admissible integrand; lowering it to $a_i \le 0$ removes the
delta function altogether, and such an integral vanishes, having no
discontinuity in that channel.  The methods developed for loop
integrals thereby become available for phase-space integration.

The partonic process behind the channel is
\begin{equation}
  p_1 + p_2 \;\longrightarrow\; p_3 + k_1 + k_2 ,
\end{equation}
with $p_1, p_2$ the incoming partons, $p_3$ the observed one, and
$k_1, k_2$ the two radiated partons that are integrated over, all
massless.  Its invariants are
\begin{equation}
  s = (p_1+p_2)^2, \qquad t = (p_1-p_3)^2, \qquad u = (p_2-p_3)^2 ,
\end{equation}
and the integrals depend on them only through the two dimensionless
ratios
\begin{equation}
  v = -\frac{t}{s}, \qquad w = -\frac{u}{s} .
\end{equation}
The physical region follows from the kinematics alone.  Since $p_3$
is massless,
\begin{equation}
  s + t + u \;=\; (p_1+p_2-p_3)^2 \;=\; M^2 ,
\end{equation}
the squared invariant mass of the unobserved pair, so that
$1 - v - w = M^2/s \ge 0$; and $v = 2p_1\!\cdot\! p_3/s > 0$,
$w = 2p_2\!\cdot\! p_3/s > 0$.  The region is therefore the open
triangle
\begin{equation}
  0 < v, \qquad 0 < w, \qquad v + w < 1 .
  \label{eq:region}
\end{equation}
In the partonic centre-of-mass frame, with $E_3$ the energy of the
observed parton and $\theta$ the angle between $\vec p_3$ and
$\vec p_1$,
\begin{equation}
  v = \frac{E_3}{\sqrt s}\,\big(1 - \cos\theta\big),
  \qquad
  w = \frac{E_3}{\sqrt s}\,\big(1 + \cos\theta\big),
  \label{eq:vwangle}
\end{equation}
so $v + w = 2E_3/\sqrt s$ fixes the energy and $v/w$ the angle.  The
three edges of the triangle are the three degenerate limits:
$v \to 0$ and $w \to 0$, in which $p_3$ becomes collinear to $p_1$
and to $p_2$; and $v + w \to 1$, the Born configuration
$E_3 = \sqrt{s}/2$, $M^2 = 0$, in which the unobserved pair carries
no invariant mass and the radiation is unresolved --- the \emph{soft
edge}, as it will be called below.  All three edges reappear throughout: as
singular loci of the differential equations, as the limits at which
boundary conditions are read off, and as the endpoints at which the
masters turn into distributions.  Everything is dimensionally
regularized, $d = 4 - 2\eps$.

\subsection{Statement of the problem}

Reduction by integration-by-parts identities leaves a basis of $347$
\emph{master integrals}.  A master is not itself a function of
$(v,w)$ alone: an $L$-loop integral whose propagator powers sum to
$a$ has mass dimension $Ld - 2a$, and Lorentz invariance then fixes
its dependence on the overall scale,
\begin{equation}
  I_m \;=\; (-s)^{\,Ld/2 - a_m}\;\hat I_m(v,w;\eps) ,
  \label{eq:dimensionless}
\end{equation}
with $\hat I_m$ dimensionless.  The differential equations below are
equations for the $\hat I_m$, and the notation $I_m(v,w;\eps)$ used
from here on refers to them; the bubble of
Example~\ref{ex:bubble} is the case $L=1$, $a=2$, where
$(-s)^{d/2-2} = (-s)^{-\eps}$ and $\hat B = c(\eps)$ is a constant.
The reduction is exact linear algebra over the field
$\Q(v,w,\eps)$, but it says nothing about the values of the basis
elements.  The task is their \emph{analytic evaluation}: for each
master, a Laurent expansion
\begin{equation}
  I(v,w;\eps) \;=\; \sum_{k \ge k_0} \eps^{k}\, I^{(k)}(v,w)
\end{equation}
whose coefficients are explicit transcendental functions --- in the
end, $\Q$-linear combinations of iterated integrals of logarithmic
differential forms --- together with exactly known constants.

The method rests on one mathematical fact: \emph{the masters satisfy
first-order linear differential equations in $v$ and $w$}.  The
sections below derive those equations, reduce the many systems to a
short list of distinct ones, construct the canonical form, treat the
square roots the kinematics forces on some of them, and determine the
boundary data and the endpoint behaviour.

%======================================================================
\section{Differential equations for master integrals}
\label{sec:de}

\subsection{Masters and integration by parts}

Fix a set of inverse propagators $D_1,\dots,D_N$ --- including the
cut ones --- large enough that every scalar product of loop and
external momenta is a linear combination of the $D_i$ and external
invariants.  Such a set, an \emph{integral family}, defines the
integrals
\begin{equation}
  I_{a_1\dots a_N}(v,w;\eps)
  \;=\; \int \prod_{\ell=1}^{2}\frac{\dd^d k_\ell}{i\pi^{d/2}}\;
  \frac{1}{D_1^{a_1}\cdots D_N^{a_N}},
  \qquad a_i \in \Z ,
  \label{eq:familyint}
\end{equation}
where for a cut line the factor $D_i^{-a_i}$ stands for the delta
function of \eqref{eq:cutpower}, and an integral
with $a_i \le 0$ on a cut line vanishes.  The \emph{sector} of an
integral is the subset of
propagators present with positive power; deleting propagators leads
to \emph{subsectors}, and the sectors of a family form a lattice
under this deletion.  Integration-by-parts identities,
\begin{equation}
  0 \;=\; \int \prod_\ell \dd^d k_\ell\;
  \frac{\partial}{\partial k_j^\mu}
  \left[\frac{q^\mu}{D_1^{a_1}\cdots D_N^{a_N}}\right],
  \qquad q \in \{k_1,k_2,p_i\},
\end{equation}
relate integrals with shifted indices; solving them expresses every
integral of the family through a finite basis of master integrals,
distributed over the sectors of the family.

\subsection{Derivation of the differential system}

A master depends on the external momenta only through their
invariants $s_{kl} = (p_k + p_l)^2$, so the chain rule gives, for
each pair $(i,j)$ of external momenta,
\begin{equation}
  p_j^\mu\,\frac{\partial}{\partial p_i^\mu}\, I
  \;=\; \sum_{k \le l}
  \Big(p_j^\mu\,\frac{\partial s_{kl}}{\partial p_i^\mu}\Big)\,
  \frac{\partial I}{\partial s_{kl}} ,
  \label{eq:chainrule}
\end{equation}
a linear system for the derivatives $\partial I/\partial s_{kl}$
whose coefficients are again invariants.  Its matrix is invertible
wherever the Gram determinant of the independent external momenta is
non-zero, and inverting it expresses each $\partial/\partial s_{kl}$
--- hence $\partial_v$ and $\partial_w$ --- as an explicit
combination of the operators $p_j^\mu\,\partial/\partial p_i^\mu$.
Those may be taken under the integral sign, where they raise
propagator powers and generate numerators $p_j\!\cdot\! k_\ell$ and
$p_i\!\cdot\! p_j$; by the completeness assumed of the family these
numerators are again linear combinations of the $D_i$ and of
invariants.  Every resulting term is thus an integral \emph{of the
same family}, and integration by parts maps it back to the masters.
Collecting
the masters of a family and its subsectors into a vector $\vec I$,
one obtains a closed linear system
\begin{equation}
  \partial_v \vec I = A_v(v,w;\eps)\,\vec I,
  \qquad
  \partial_w \vec I = A_w(v,w;\eps)\,\vec I,
  \label{eq:desys}
\end{equation}
with matrices of rational functions.  It is convenient to combine the
pair into the matrix-valued one-form $A = A_v\,\dd v + A_w\,\dd w$,
the \emph{connection}, and to write
\begin{equation}
  \dd \vec I \;=\; A\,\vec I .
  \label{eq:pfaffian}
\end{equation}
A total differential equation of this shape --- one first-order
equation for the differential of an unknown vector, in several
variables at once --- is what the literature calls a \emph{Pfaffian
system}, and \eqref{eq:pfaffian} will be referred to as the Pfaffian
system of the family.  A Pfaffian system is not solvable for arbitrary
$A$: the two equations of \eqref{eq:desys} constrain each other, and
the condition under which they are compatible is derived in
Section~\ref{sec:integrability}.  The differential-equation method for
master integrals goes back to Kotikov~\cite{kotikov} and to Gehrmann
and Remiddi~\cite{gr}.

\begin{example}[The one-mass bubble]
\label{ex:bubble}
Let
\begin{equation}
  B(s;\eps) = \int \frac{\dd^d k}{i\pi^{d/2}}
  \frac{1}{k^2\,(k+p)^2}, \qquad s = p^2 ,
\end{equation}
a one-scale family with $D_1 = k^2$, $D_2 = (k+p)^2$ and a single
nontrivial master, $B$ itself.  Here \eqref{eq:chainrule} has a
single term, $p^\mu\partial_{p^\mu} = 2s\,\partial_s$, so that
$\partial_s = \frac{1}{2s}\, p^\mu \partial_{p^\mu}$ and
\begin{equation}
  \frac{\partial B}{\partial s}
  = -\frac{1}{2s}\int \frac{\dd^d k}{i\pi^{d/2}}
    \frac{2p\cdot(k+p)}{k^2\,[(k+p)^2]^2} .
\end{equation}
The numerator is a combination of inverse propagators,
$2 p\cdot(k+p) = D_2 - D_1 + s$, so the right-hand side involves
only integrals of the family: $I_{1,1} = B$, the scaleless
$I_{0,2} = 0$, and $I_{1,2}$.  The integration-by-parts identity
generated by $\partial_{k^\mu}\!\big(k^\mu/D_1 D_2\big)$ reads
$0 = (d-3)\,I_{1,1} + s\,I_{1,2}$, that is,
\begin{equation}
  I_{1,2} = -\frac{d-3}{s}\, I_{1,1},
\end{equation}
and therefore
\begin{equation}
  \frac{\partial B}{\partial s}
  = -\frac{1}{2s}\Big[\,B - (d-3)\,B\,\Big]
  = -\,\frac{4-d}{2s}\,B
  = -\,\frac{\eps}{s}\, B .
  \label{eq:bubblede}
\end{equation}
The equation determines the entire $s$-dependence,
\begin{equation}
  B(s;\eps) \;=\; c(\eps)\,(-s)^{-\eps},
\end{equation}
with the branch fixed by the causal prescription $s \to s + i0$ ---
consistent with dimensional analysis, since $B$ carries mass
dimension $d - 4 = -2\eps$.  A single direct evaluation fixes the
remaining constant,
\begin{equation}
  c(\eps) \;=\;
  \frac{\Gamma(1+\eps)\,\Gamma(1-\eps)^2}
       {\eps\,(1-2\eps)\,\Gamma(1-2\eps)} .
\end{equation}
The pattern is the one followed throughout: the differential equation
fixes the functional dependence, and a single elementary evaluation
fixes the constant left free by it.  The value of $c(\eps)$ is
recovered from the parametric representation in
Section~\ref{sec:classes}.
\end{example}

\subsection{Triangular structure}

Order the sectors of a family by their number of propagators.
Differentiating an integral of a sector $S$ raises the powers of the
lines already present; the numerators so produced are rewritten
through the $D_i$, which may cancel a line and land in a subsector of
$S$, as it does for $I_{0,2}$ in Example~\ref{ex:bubble}.  The
integration-by-parts identities likewise express any integral of $S$
through integrals of $S$ and of its subsectors, since an identity may
cancel a line but never create one.  Neither operation can produce a
line that was absent, so both lead sideways and down the sector
lattice, never up it.  With the master
vector ordered accordingly, the connection is
\emph{block lower triangular},
\begin{equation}
  A \;=\;
  \begin{pmatrix}
    a_{11} & & & \\
    a_{21} & a_{22} & & \\
    \vdots & & \ddots & \\
    a_{b1} & a_{b2} & \cdots & a_{bb}
  \end{pmatrix},
  \label{eq:blocktriangular}
\end{equation}
where each diagonal entry $a_{ss}$ --- a \emph{block}, in the
terminology fixed in the next section --- couples the masters of one
group of sectors among themselves, and the entries below the
diagonal express how lower sectors source higher ones.  The diagonal
blocks carry the irreducible analytic content of the system; the
entries below the diagonal can be dealt with afterwards --- the
observation on which Section~\ref{sec:finding} rests.

\subsection{Integrability}
\label{sec:integrability}

Equality of mixed partial derivatives,
$\partial_v\partial_w \vec I = \partial_w\partial_v \vec I$, forces
the compatibility condition
\begin{equation}
  \partial_v A_w - \partial_w A_v - [A_v, A_w] = 0,
  \qquad\text{equivalently}\qquad
  \dd A \;=\; A \wedge A ,
  \label{eq:flatness}
\end{equation}
the statement that the connection is \emph{flat}.  This is the
integrability condition of the Pfaffian system \eqref{eq:pfaffian}:
it is what makes the two equations of \eqref{eq:desys} compatible,
and a Pfaffian system satisfying it is called integrable.  (To check
the sign: $\partial_v\partial_w \vec I = \partial_v(A_w \vec I) =
(\partial_v A_w)\vec I + A_w A_v \vec I$, and antisymmetrizing in
$v,w$ gives $\partial_v A_w - \partial_w A_v + [A_w, A_v] = 0$.)
Flatness is a stringent consistency condition on any derived or
transformed system, since it interlocks all matrix entries; it is
preserved by every change of basis; and it guarantees that solutions
continued along a path depend only on the endpoint and on the
homotopy class of the path in the complement of the singular loci,
which underlies Section~\ref{sec:paths}.

%======================================================================
\section{Families, blocks, and equivalence}
\label{sec:classes}

The double-real channel does not present one large differential
system but many small ones, and many of the small ones coincide as
differential systems once their bases are matched.  Two equivalence
relations --- one among integral families, one among diagonal blocks
--- detect the coincidences and together set the true size of the
analytic problem.  Both are decided by exact identities: a change of
loop momenta for families, a conjugation of matrices for blocks.

\subsection{Equivalence of integral families}

Families are generated mechanically, one propagator set per diagram,
and diagrams proliferate: the channel begins with $1898$ of them.
Most are redundant, because different routings of the loop momenta
describe the same set of denominators.

\begin{definition}[Equivalent families; representative families]
\label{def:famequiv}
Two integral families with $L$ loop momenta are \emph{equivalent} if
an \emph{affine} change of the loop momenta --- a linear map composed
with a shift, here the matrix $M$ composed with a shift by external
momenta ---
\begin{equation}
  k_\ell \;\longmapsto\; \sum_{m=1}^{L} M_{\ell m}\, k_m
  \;+\; \sum_{r} c_{\ell r}\, p_r ,
  \qquad M \in \mathrm{GL}(L,\Z),\quad c_{\ell r} \in \Q ,
  \label{eq:shift}
\end{equation}
maps one propagator set bijectively onto the other and maps cut
lines to cut lines.  A matrix in $\mathrm{GL}(L,\Z)$ has
$\lvert\det M\rvert = 1$, so the integration measure is invariant and
equivalent families have identical integrals up to a relabeling of
the indices $a_i$.  A chosen representative of each equivalence class
is called a \emph{representative family}.
\end{definition}

\begin{example}[Two routings of the one-mass triangle]
\label{ex:shift}
Let $p_1^2 = p_2^2 = 0$ and $(p_1+p_2)^2 = s$, and consider the two
one-loop families
\begin{equation}
  \mathcal{F}:\ \{\,k^2,\ (k+p_1)^2,\ (k+p_1+p_2)^2\,\},
  \qquad
  \widetilde{\mathcal{F}}:\ \{\,k^2,\ (k+p_1+p_2)^2,\ (k+p_2)^2\,\},
\end{equation}
the same triangle graph with its external legs attached in a
different order.  No relabeling of the $D_i$ maps one set onto the
other.  But the substitution $k \mapsto -k - p_1 - p_2$, of the
form \eqref{eq:shift} with $M = (-1)$, maps
$\widetilde{\mathcal{F}}$ onto $\mathcal{F}$:
\begin{equation}
  k^2 \mapsto (k+p_1+p_2)^2, \qquad
  (k+p_1+p_2)^2 \mapsto k^2, \qquad
  (k+p_2)^2 \mapsto (k+p_1)^2 ,
\end{equation}
a bijection of the propagator sets.  The families are equivalent,
and every integral of one is an integral of the other with permuted
indices.
\end{example}

Searching over the pairs $(M, c)$ directly is not how equivalence is
found: the search space is large, and the failure of a search proves
nothing.  Instead one computes an invariant that does not depend on
the routing at all.  For positive powers, the integral
\eqref{eq:familyint} has the Feynman-parametric representation
\begin{equation}
  I_{a_1\dots a_N}
  \;=\; (-1)^{a}\,
  \frac{\Gamma\!\big(a - L d/2\big)}{\prod_i \Gamma(a_i)}
  \int_{x_i \ge 0} \Big(\prod_i \dd x_i\, x_i^{a_i-1}\Big)\,
  \delta\Big(1 - \sum_i x_i\Big)\,
  \frac{U^{\,a-(L+1)d/2}}{F^{\,a-Ld/2}} ,
  \qquad a = \sum_i a_i ,
  \label{eq:parametric}
\end{equation}
where $U$ and $F$ are the \emph{Symanzik polynomials} of the graph,
\begin{equation}
  U \;=\; \sum_{T} \prod_{e \notin T} x_e ,
  \qquad
  F \;=\; -\sum_{(T_1,T_2)} s_{(T_1,T_2)}
  \prod_{e \notin T_1 \cup T_2} x_e
  \qquad \text{(massless lines)} :
  \label{eq:symanzik}
\end{equation}
$T$ runs over the spanning trees of the graph, $(T_1,T_2)$ over its
spanning two-forests, and $s_{(T_1,T_2)}$ is the square of the total
external momentum flowing into $T_1$.  At one loop this reduces to
$F = -\sum_{i<j} q_{ij}^2\, x_i x_j$ with $q_{ij}$ the difference of
the momenta carried by lines $i$ and $j$.  For the bubble of
Example~\ref{ex:bubble} ($L = 1$, $a_1 = a_2 = 1$, $U = x_1 + x_2$,
$F = -s\,x_1x_2$), the representation \eqref{eq:parametric} restricted
to the simplex $x_1 + x_2 = 1$ gives
\begin{equation}
  B \;=\; \Gamma(\eps)\,(-s)^{-\eps}
  \int_0^1 \dd x\,\big[x(1-x)\big]^{-\eps}
  \;=\; \frac{\Gamma(\eps)\,\Gamma(1-\eps)^2}{\Gamma(2-2\eps)}\,
  (-s)^{-\eps} \;=\; c(\eps)\,(-s)^{-\eps} ,
  \label{eq:bubbleparametric}
\end{equation}
using $\Gamma(1+\eps) = \eps\,\Gamma(\eps)$ and
$\Gamma(2-2\eps) = (1-2\eps)\Gamma(1-2\eps)$ in the last step: the
normalization of \eqref{eq:parametric}, including its sign, and the
constant of Example~\ref{ex:bubble} confirm each other.  The essential
point is
visible in \eqref{eq:symanzik}: $U$ and $F$ are built from the graph
and from the assignment of external momenta to its vertices --- the
loop momenta have disappeared.  A substitution \eqref{eq:shift}
therefore acts on the parametric representation as a mere
permutation of the Feynman parameters, and so
\emph{equivalent families have identical Symanzik pairs $(U,F)$ up
to a permutation of the parameters}.

\begin{example}[The triangle, in parameters]
\label{ex:pak}
For the two routings of Example~\ref{ex:shift},
\begin{equation}
  U = x_1 + x_2 + x_3 \ \ \text{for both},
  \qquad
  F_{\mathcal{F}} = -s\, x_1 x_3 ,
  \qquad
  F_{\widetilde{\mathcal{F}}} = -s\, x_1 x_2 ,
\end{equation}
by \eqref{eq:symanzik} or by the one-loop formula ($q_{13} =
p_1{+}p_2$ for $\mathcal{F}$, $q_{12} = p_1{+}p_2$ for
$\widetilde{\mathcal{F}}$, all other $q_{ij}^2$ vanishing): the
pairs coincide after $x_2 \leftrightarrow x_3$.  The shift of
Example~\ref{ex:shift} acts on the parameters as the relabeling
$\tilde x_1 \mapsto x_3$, $\tilde x_2 \mapsto x_1$,
$\tilde x_3 \mapsto x_2$ dictated by its action on the propagators,
under which $-s\,\tilde x_1\tilde x_2 \mapsto -s\,x_3x_1$ as well.
The two matchings differ by $x_1 \leftrightarrow x_3$, a symmetry of
the pair $(x_1{+}x_2{+}x_3,\ -s\,x_1x_3)$; matchings of the
polynomials always form a coset of that symmetry group, and each
member is a candidate for the underlying momentum shift.  They are
found on the polynomials, without any search over momenta.
\end{example}

Comparing polynomial pairs up to permutation is itself a finite
problem, and Pak's algorithm~\cite{pak} settles it canonically.  Write
the pair
as a matrix with one row per monomial: the entries of a row are the
exponents of the parameters in that monomial, followed by the
coefficient, with the rows of $U$ and of $F$ kept in separate groups.
For the triangle of Example~\ref{ex:pak} the matrix is
\begin{equation}
  \begin{array}{c|ccc|c}
     & x_1 & x_2 & x_3 & \text{coeff.}\\\hline
   U & 1 & 0 & 0 & 1\\
   U & 0 & 1 & 0 & 1\\
   U & 0 & 0 & 1 & 1\\\hline
   F & 1 & 0 & 1 & -s
  \end{array}
  \qquad\text{for } \mathcal{F},
  \qquad
  \text{and the row } (1,1,0) \text{ in place of } (1,0,1)
  \text{ for } \widetilde{\mathcal{F}} .
\end{equation}
A relabeling of the parameters permutes the columns; the rows are then
sorted, and one seeks the relabeling making the sorted matrix
lexicographically maximal.  The search is greedy and refines column by
column: at step $j$ each surviving partial relabeling is extended by
each parameter not yet placed, the resulting matrices are compared on
their first $j$ columns, and only the maximizers survive; the
relabelings that survive to the end form a coset of the symmetry group
of the pair, and are that group when the pair is already normally
ordered.  Here $U$ is symmetric and only the $F$ row decides: among
$(1,1,0)$, $(1,0,1)$ and $(0,1,1)$ the maximum is $(1,1,0)$, so both
routings normalize to
\begin{equation}
  \big(x_1+x_2+x_3,\ -s\,x_1x_2\big),
\end{equation}
reached from $\widetilde{\mathcal{F}}$ by the identity and from
$\mathcal{F}$ by $x_2 \leftrightarrow x_3$, with the residual
symmetry $x_1 \leftrightarrow x_2$ left over.  The outcome is a
\emph{normal ordering} of the parameters, the same for every member
of a permutation orbit, so two families are candidates for
equivalence exactly when their normally ordered pairs are identical.
Cut lines are carried through with distinguishing marks, so that the
normal ordering can never trade a cut line for an uncut one; for a
cut family the parametric representation serves only as this
fingerprint of the underlying graph, never for evaluation.

A coincidence of normal orderings is a reliable \emph{candidate} for
equivalence; every candidate is then confirmed by constructing the
substitution \eqref{eq:shift} explicitly.  The permutation $\sigma$
of parameters dictates which propagator must map onto which, and the
requirement
\begin{equation}
  \widetilde D_i\big(Mk + c\cdot p\big) \;=\; D_{\sigma(i)}(k)
  \qquad \text{for every } i
  \label{eq:shiftconditions}
\end{equation}
is solved as follows: writing each inverse propagator as the square of
a linear form, \eqref{eq:shiftconditions} equates the two linear forms
up to sign, one vector equation per line.  The system is linear in
$(M,c)$ once a sign has been chosen for each line; the sign choices
are a finite branching, and at one loop they are trivial, the sign
being absorbed into $M$.  There, with $D_j = (k+r_j)^2$ and
$\widetilde D_i = (k+\tilde r_i)^2$, the conditions read
$M\,\big(c\!\cdot\!p + \tilde r_i\big) = r_{\sigma(i)}$; for
Example~\ref{ex:shift}, $M = -1$ and $c\!\cdot\!p = -(p_1{+}p_2)$
solve $p_1{+}p_2-\tilde r_i = r_{\sigma(i)}$ for the three lines at
once.
Applying the substitution so obtained is what turns the identification
into a proof.  The $1898$ raw families of the channel collapse to
$430$ inequivalent ones, and only a fraction of the survivors
supports master integrals at all: ninety-one families, in the end,
carry the Pfaffian systems of the problem.

Differentiation is family-local --- the derivative of a master of a
family reduces to masters of that family and its subsectors, never
to unrelated ones --- so the differential-equation problem separates
into $91$ Pfaffian systems, of dimensions between $2$ and $47$,
accounting for all $347$ masters.  (The dimensions sum to more than $347$: each
system contains, besides its own masters, the subsector masters it
couples to, and those are counted once, in the family that owns
them.)  The systems are separate but not independent: a subsector
master shared by two of them is one and the same function, so they
are solved in an order compatible with the subsector relation, as
Section~\ref{sec:boundary} uses.  The boundaries of the systems were
fixed by the reduction itself; what must be constructed anew is, for
each family, the connection $(A_v, A_w)$, together with the proof of
its flatness.

\subsection{Blocks and their equivalence}

\begin{definition}[Block]
Form the directed graph whose vertices are the masters of a system
and which has an edge $i \to j$ whenever $A_{ji} \not\equiv 0$, that
is, whenever $I_i$ occurs in the derivative of $I_j$.  A \emph{block}
is a strongly connected component of this graph: a set of masters in
which every two are joined in both directions by chains of such
couplings.  With the masters ordered by sector, the blocks are the
diagonal entries $a_{ss}$ of Eq.~\eqref{eq:blocktriangular}, and a
block of size one is a master whose derivative involves only itself
and lower sectors.
\end{definition}

\begin{definition}[Equivalence of blocks; classes]
\label{def:class}
Let $b^{(1)}$ and $b^{(2)}$ be the connections of two blocks, in the
same family or in different ones.  The blocks are \emph{equivalent}
if there is a permutation matrix $P$ with
\begin{equation}
  b^{(2)}(v,w) \;=\; P^{-1}\, b^{(1)}(v,w)\, P
  \qquad\text{or}\qquad
  b^{(2)}(v,w) \;=\; P^{-1}\, b^{(1)}(w,v)\, P ,
  \label{eq:blockequiv}
\end{equation}
as an identity of rational matrices.  An equivalence class of
blocks is called a \emph{class}; each class is represented by a
single block, whose canonical form is constructed once and serves
for all members.
\end{definition}

Because $P$ is constant, the conjugation in \eqref{eq:blockequiv}
carries no derivative term ($P^{-1}\dd P = 0$): the equivalence is a
plain identity of rational matrices, decided entry by entry.  In the
double-real channel the $91$ family Pfaffian systems contain $1119$
blocks,
and the $1119$ blocks fall into $173$ classes.

\begin{example}[Two equivalent blocks]
\label{ex:blockequiv}
Suppose one family contains the $2\times2$ block
\begin{equation}
  b^{(1)} = \eps\left[
  \begin{pmatrix} 0 & 0\\ 0 & -1\end{pmatrix}\dd\log w
  +
  \begin{pmatrix} 0 & 0\\ 1 & 0 \end{pmatrix}\dd\log(1-w)
  \right]
\end{equation}
and another family the block
\begin{equation}
  b^{(2)} = \eps\left[
  \begin{pmatrix} -1 & 0\\ 0 & 0\end{pmatrix}\dd\log v
  +
  \begin{pmatrix} 0 & 1\\ 0 & 0 \end{pmatrix}\dd\log(1-v)
  \right] .
\end{equation}
With $P = \begin{psmallmatrix} 0 & 1\\ 1 & 0\end{psmallmatrix}$
(reversal of the basis order) and the swap $v \leftrightarrow w$,
conjugation sends
$\begin{psmallmatrix} 0 & 0\\ 0 & -1\end{psmallmatrix} \mapsto
 \begin{psmallmatrix} -1 & 0\\ 0 & 0\end{psmallmatrix}$ and
$\begin{psmallmatrix} 0 & 0\\ 1 & 0\end{psmallmatrix} \mapsto
 \begin{psmallmatrix} 0 & 1\\ 0 & 0\end{psmallmatrix}$,
so the second identity in \eqref{eq:blockequiv} holds, and any
canonical solution of $b^{(1)}$ becomes one of $b^{(2)}$ after the
swap and the reversal.  The block $b^{(1)}$ is the system whose
solution
will be constructed in Example~\ref{ex:li2}, where its second
component turns out to be the dilogarithm.
\end{example}

The roughly sixfold collapse from blocks to classes has a physical
origin.  A block is the differential system of the masters of one
subtopology, and the same subtopology occurs as a subdiagram of
many different parent diagrams: the identification of
\emph{complete} families by momentum shifts leaves untouched the far
more frequent coincidence of their \emph{subsectors}.  The
involution $v \leftrightarrow w$ is the exchange
$p_1 \leftrightarrow p_2$ of the incoming partons, that is
$t \leftrightarrow u$ --- a relabeling of the initial state, not a
crossing --- which maps cut diagrams into cut diagrams.  Equivalence
is never assumed on such grounds, however: a candidate identification
is accepted only when \eqref{eq:blockequiv} is verified identically.

The analysis of Sections~\ref{sec:finding} and \ref{sec:roots} is
therefore carried out once per class rather than once per block, and
each family is afterwards reconstructed from the canonical forms of
its classes.  Which rows of a family realize which class, and under
which permutation of the basis and which exchange of the variables, is
part of the statement being proved and is verified again at the
reconstruction.

%======================================================================
\section{The canonical form}
\label{sec:epsform}

\subsection{The $\eps$-form and its immediate consequences}

\begin{definition}[Canonical form; letters, alphabet, residues]
A Pfaffian system is in \emph{canonical form} (equivalently, in
\emph{$\eps$-form}; both names are standard, after
Henn~\cite{henn}) if its connection is
\begin{equation}
  \dd \vec F \;=\; \eps \sum_{a} R_a \,\dd\log \varphi_a\;
  \vec F ,
  \label{eq:epsform}
\end{equation}
where the \emph{residue matrices} $R_a$ are constant --- independent
of kinematics and regulator alike --- and the \emph{letters}
$\varphi_a$ are functions of the kinematics alone.  The collection
$\{\varphi_a\}$ is the \emph{alphabet} of the system.  In the
simplest and most common case the letters are polynomials;
Section~\ref{sec:roots} enlarges the notion to algebraic functions.
\end{definition}

One family of the double-real channel, a system of dimension
fourteen, carries the six-letter alphabet
\begin{equation}
  \{\, v,\; w,\; 1-v,\; 1-w,\; v+w,\; 1-v-w \,\} ,
\end{equation}
and every entry of its canonical connection is $\eps$ times a
$\Q$-linear combination of logarithmic differentials of these six
polynomials.  Richer systems have richer alphabets --- among the
two-root families one finds letters quadratic in the variables ---
but always of the same form.

The equivalence of blocks, Definition~\ref{def:class}, can now be
stated on the data of \eqref{eq:epsform}: for two blocks in canonical
form the identity \eqref{eq:blockequiv} reduces to a statement about
the residue matrices alone,
\begin{equation}
  R^{(2)}_a \;=\; P^{-1} R^{(1)}_{\sigma(a)} P ,
  \label{eq:classresidues}
\end{equation}
where $\sigma$ is the relabeling of the letters induced by the
exchange (under $v \leftrightarrow w$ the letter $v$ becomes $w$,
while $1-v-w$ is fixed, and so on).  A class is thus a finite set of
constant matrices up to simultaneous conjugation and relabeling.

Three consequences follow at once, and they are linked.  The
connection has simple poles exactly on the loci $\varphi_a = 0$ and
nowhere else in the finite plane --- the system is \emph{Fuchsian},
and for generic $\eps$ each letter with non-vanishing residue is a
genuine branch locus of the solutions.  (Infinity generally carries a
pole as well; Section~\ref{sec:diagonal} makes use of it.)  Every
solution is therefore analytic on the complement of
$\bigcup_a \{\varphi_a = 0\}$ and has all its finite branch points
among the zeros of letters: the alphabet is the complete list of
singular loci.  Second, because each $\dd\log\varphi_a$ is a closed
form, the flatness condition \eqref{eq:flatness} becomes a condition
on the commutators of the constant matrices, weighted by rational
functions,
\begin{equation}
  \sum_{a<b} [R_a, R_b]\; W_{ab}(v,w) \;=\; 0 ,
  \qquad
  W_{ab} \;=\;
  \frac{\partial_v\varphi_a\,\partial_w\varphi_b
      - \partial_w\varphi_a\,\partial_v\varphi_b}
       {\varphi_a\,\varphi_b} ,
  \label{eq:flatletters}
\end{equation}
where $\dd\log\varphi_a \wedge \dd\log\varphi_b =
W_{ab}\,\dd v \wedge \dd w$.  The $W_{ab}$ are rational, not
independent, and the independent conditions on the commutators are
obtained by decomposing them in partial fractions and equating the
coefficient of each independent pole structure to zero; that done,
flatness is a finite list of commutator identities, checked once and
for all.  Third, as the next subsection shows, the coefficient of
$\eps^k$ in the solution is a $\Q$-linear combination of $k$-fold
iterated integrals contracted with the boundary vector, so the
expansion is homogeneous in \emph{transcendental weight} --- the
grading on functions and numbers under which a logarithm carries
weight one, a $k$-fold iterated integral of logarithmic forms weight
$k$, and the constants $\zeta_k$ and $\pi^k$ weight $k$, with
rational numbers at weight zero.  In the canonical basis this grading
is manifest, order by order in $\eps$.

\subsection{Order-by-order integration}

In a generic basis the regulator is entangled with the kinematics in
every entry of $A$, and the coefficients of the $\eps$-expansion obey
coupled differential equations with no distinguished structure.  In
the canonical basis the regulator counts the number of integrations
and does nothing else, and the system reduces to quadratures.

To see this, write $\hat A = \sum_a R_a\,\dd\log\varphi_a$, so that
$\dd\vec F = \eps\,\hat A\,\vec F$ with $\hat A$ independent of
$\eps$.  Substituting $\vec F = \sum_{k}\eps^k \vec F^{(k)}$ and
comparing powers of $\eps$:
\begin{equation}
  \dd \vec F^{(0)} = 0,
  \qquad
  \dd \vec F^{(k)} = \hat A\, \vec F^{(k-1)} \qquad (k \ge 1) .
  \label{eq:orderbyorder}
\end{equation}
The system has disentangled into a recursion of quadratures: the
lowest order is a constant vector, and each subsequent order is an
integral of known functions.  Along a path $\gamma$ from a base
point $p_0$,
\begin{equation}
  \vec F^{(k)}(p) = \vec F^{(k)}(p_0)
  + \int_\gamma \hat A\;\vec F^{(k-1)} ,
\end{equation}
and unwinding the recursion expresses the full solution as the
path-ordered exponential
\begin{equation}
  \vec F(p) \;=\; \Pexp\!\Big(\eps \int_{\gamma} \hat A\Big)\,
  \vec F(p_0)
  \;=\; \sum_{k\ge0} \eps^k \!\!
  \sum_{a_1,\dots,a_k}\!\!
  R_{a_k}\!\cdots R_{a_1}
  \int\limits_{0\le t_1\le\cdots\le t_k\le 1}
  \prod_{i=1}^{k}\dd\log \varphi_{a_i}\big(\gamma(t_i)\big)\;
  \vec F(p_0) .
  \label{eq:pexp}
\end{equation}
The $k$-fold integrals are \emph{Chen iterated
integrals}~\cite{chen} over the alphabet; flatness makes them
independent of the choice of path within its homotopy class.  The analytic problem is thereby solved
in principle, and what remains is to express the iterated integrals
in a standard basis of functions (Section~\ref{sec:paths}) and to
determine the boundary vector $\vec F(p_0)$
(Section~\ref{sec:boundary}).

The recursion \eqref{eq:orderbyorder} starts from a \emph{constant}
at order $\eps^0$, while the masters themselves have poles in $\eps$.
There is no contradiction: the matrix $T$ relating the derived
masters to the canonical ones carries inverse powers of $\eps$ from
the normalization of the basis, so a finite canonical order feeds
several Laurent orders of a master, and statements about uniform
weight refer to the canonical basis.

In one dimension $\dd F = \eps\,\dd\log w\,F$ has the solution
$F = C\,w^{\eps} = C\big(1 + \eps\log w + \tfrac{\eps^2}{2}\log^2 w
+ \cdots\big)$, and the recursion \eqref{eq:orderbyorder} generates
exactly that expansion.  Both descriptions of this object --- the
expanded logarithmic series and the unexpanded power $w^\eps$ ---
will be needed, and the tension between them is the subject of
Section~\ref{sec:endpoint}.

\begin{example}[A two-dimensional system and the dilogarithm]
\label{ex:li2}
Take the alphabet $\{w,\,1-w\}$ on $0<w<1$ and the system
\begin{equation}
  \dd \vec F = \eps\left[
  \begin{pmatrix} 0 & 0\\ 0 & -1\end{pmatrix}\dd\log w
  +
  \begin{pmatrix} 0 & 0\\ 1 & 0 \end{pmatrix}\dd\log(1-w)
  \right]\vec F ,
\end{equation}
with the boundary condition $\vec F \to (1,0)^{\mathsf T}$ as
$w \to 0$ (the precise sense in which a value is assigned at a
singular point is given in Section~\ref{sec:local}).  The first
component stays constant.  For the second, the recursion gives
\begin{align}
  F_2^{(1)}(w) &= \int_0^w \dd\log(1-w')\cdot 1
               \;=\; \log(1-w) ,\\
  F_2^{(2)}(w) &= \int_0^w \Big[\dd\log(1-w')\cdot 0
               \;-\;\dd\log w' \cdot F_2^{(1)}(w')\Big]
               = -\int_0^w \frac{\dd w'}{w'}\,\log(1-w')
               \;=\; \operatorname{Li}_2(w) .
\end{align}
The dilogarithm appears as the twofold iterated integral with the
word $(w,\,1{-}w)$; the next order produces trilogarithms, and so
on.  Every master of the calculation is, order by order, a
$\Q$-linear combination of such words in its own alphabet.
\end{example}

\subsection{Local behaviour at the singular loci}
\label{sec:local}

The canonical form does more than integrate: it displays, before any
integration, how every solution behaves at every singular locus.
Near a zero of a letter --- say $w \to 0$, with letter $w$ --- the
system reads
\begin{equation}
  \partial_w \vec F \;=\; \frac{\eps\,R_{w}}{w}\,\vec F
  \;+\; (\text{regular}) ,
\end{equation}
and the classical Frobenius analysis applies: for generic $\eps$ every
solution is
\begin{equation}
  \vec F \;=\; H(w)\;w^{\,\eps R_w}\,\vec c
  \;=\; \sum_{\text{modes}} w^{\,n + a\eps}\,\big(\log w\big)^{j}\,
  \big(\text{const} + O(w)\big) ,
  \label{eq:frobenius}
\end{equation}
with $H$ analytic at the origin and $H(0) = 1$, $\vec c$ a constant
vector, and the powers $w^{\eps R_w} = e^{\eps R_w \log w}$ giving the
modes on the right: the coefficients $a$ range over the eigenvalues of
$R_w$,
non-negative integer shifts $n$ from the Taylor expansion of the
regular part, and powers of $\log w$ from the Jordan structure of
$R_w$.  \emph{The complete local behaviour of every solution at
every singular locus is thus encoded in the residue matrix at the
corresponding letter} --- exponents in its eigenvalues, logarithms
in its Jordan blocks --- and is known exactly in advance.  Matching
the global solution onto \eqref{eq:frobenius} is also what assigns
boundary data at a singular point: the data are the components of
$\vec c$, that is the coefficients of the modes in the Frobenius
decomposition, of which the coefficient of the mode with exponent
zero is the ordinary value when the solution is regular there.

Two qualifications keep the statement honest.  The normal form
requires that no two exponents $\eps\lambda_i$, $\eps\lambda_j$
differ by a nonzero integer --- automatic for generic $\eps$, which
is the regime of the Laurent expansion, but not an identity in
$\eps$.  And ``logarithms come from the Jordan blocks'' refers to
the solution at fixed $\eps$; upon expanding in $\eps$ the
logarithms proliferate for a second, independent reason,
\begin{equation}
  w^{\,\eps R_w} \;=\; e^{\,\eps R_w \log w}
  \;=\; \sum_{m \ge 0} \frac{(\eps \log w)^m}{m!}\, R_w^{\,m} ,
\end{equation}
so a power of $\log w$ appears at every order in $\eps$ whether or
not $R_w$ is diagonalizable; the Jordan structure bounds only the
logarithms that survive at fixed $\eps$.  This distinction --- the
unexpanded mode versus its logarithmic series --- is precisely what
the endpoint analysis of Section~\ref{sec:endpoint} turns on.

\begin{remark}[Existence]
No theorem guarantees that a given block admits a canonical form with
rational letters, and a system with an irregular singularity admits
none at all (Example~\ref{ex:fuchs}).  The canonical form is
accordingly treated as a \emph{statement to be verified}: any means
whatsoever may produce the transformation, and what establishes it is
the exact identity between the transformed connection and the form
\eqref{eq:epsform}.
\end{remark}

In the double-real channel every class admits a canonical form, in a
number of cases only after a change of variables that rationalizes a
square root; for the families whose alphabets carry three independent
roots no such change of variables exists, and
Section~\ref{sec:roots} describes what takes its place.

%======================================================================
\section{Construction of the canonical transformation}
\label{sec:finding}

The object to be constructed is an invertible matrix $T(v,w;\eps)$
relating the derived basis to the canonical one,
$\vec I = T\,\vec F$.  Substituting into the Pfaffian system
\eqref{eq:pfaffian} shows that the connection transforms
inhomogeneously,
\begin{equation}
  A \;\longmapsto\; A' \;=\; T^{-1} A\, T \;-\; T^{-1}\dd T ,
  \label{eq:gauge}
\end{equation}
and the requirement is that $A'$ be of the form \eqref{eq:epsform}.
The triangular structure \eqref{eq:blocktriangular} decides the order
in which the blocks are treated: a diagonal block never feels what
lies below it, so each is brought to canonical form on its own, once
for each class, after which the entries below the diagonal yield to
transformations of unipotent shape.  What survives both steps is a
dependence on $\eps$ constant in the kinematics, removed last by
conjugation.

\subsection{Diagonal blocks: Fuchsian reduction and balances}
\label{sec:diagonal}

For the diagonal blocks two complementary strategies are classical,
and each produces a candidate transformation which then stands or
falls on the exact identity.

The first is the \emph{rational ansatz}.  Multiplying
\eqref{eq:gauge} by $T$ from the left removes the inverse and leaves
an identity between rational one-forms,
\begin{equation}
  \dd T \;=\; A\,T \;-\; \eps\, T \sum_a R_a\,\dd\log\varphi_a .
  \label{eq:ansatzeq}
\end{equation}
The candidate alphabet is not guessed: the letters must account for
the singularities the system already has, so one takes the
irreducible factors of the denominators of $A$, together with the
irreducible factors of the discriminants that enter the eigenvalues
of its residue matrices --- the latter being where algebraic letters,
if any, come from.  One then posits for $T$ a matrix of rational
functions whose denominators are products of these letters and whose
numerators are polynomials of bounded degree with undetermined
$\eps$-dependent coefficients.  Clearing denominators and comparing
coefficients of the independent monomials turns \eqref{eq:ansatzeq}
into a finite algebraic system in the unknowns.  It is bilinear ---
the products $T R_a$ pair the numerator coefficients with the
residues --- so it is not one linear solve: one fixes the residues
from a finite list of candidates and solves linearly for $T$ in each
case, or eliminates the residues first.  With that caveat the ansatz
is exhaustive at a given degree and a given alphabet: if the system
has no solution, no transformation of that degree exists over that
alphabet.

The second strategy proceeds through normal forms, in two steps: a
higher pole is first reduced, if it can be, to a simple one, and the
residue eigenvalues are then normalized.  The first step is Moser's
algorithm~\cite{moser}: a rank condition on the leading coefficient of
the pole
decides whether a rational transformation lowering the Moser rank
exists and constructs one when it does; a pole that the criterion
cannot lower is a genuine irregular singularity, and then no
canonical form exists at all (Example~\ref{ex:fuchs}).  The second
step uses \emph{balance transformations} between a pair of singular
points $v_1, v_2$,
\begin{equation}
  T \;=\; \bar P + \frac{v - v_1}{\,v - v_2\,}\, P ,
  \qquad \bar P = 1 - P ,
  \label{eq:balance}
\end{equation}
with $P$ a projector of rank one.  Which projector is not free.
Writing $A_{(1)}, A_{(2)}$ for the residues at $v_1, v_2$ and using
$T^{-1} = \bar P + f^{-1}P$ with $f = (v-v_1)/(v-v_2)$,
\begin{equation}
  T^{-1}A\,T \;=\;
  \bar P A \bar P \;+\; f\,\bar P A P
  \;+\; f^{-1} P A \bar P \;+\; P A P ,
\end{equation}
and the third term has a double pole at $v_1$, the second a double
pole at $v_2$, unless
\begin{equation}
  P\,A_{(1)}\,\bar P \;=\; 0
  \qquad\text{and}\qquad
  \bar P\,A_{(2)}\,P \;=\; 0 .
  \label{eq:balanceconditions}
\end{equation}
Both hold if one takes a \emph{left} eigenvector $\eta$ of $A_{(1)}$,
with eigenvalue $\lambda$, and a \emph{right} eigenvector $u$ of
$A_{(2)}$, with eigenvalue $\mu$, such that $\eta^{\mathsf T}u \ne 0$,
and sets
\begin{equation}
  P \;=\; \frac{u\,\eta^{\mathsf T}}{\eta^{\mathsf T}u} .
\end{equation}
A spectral projector of the residue at $v_1$ alone will not do: it
satisfies the first condition and in general violates the second.
For $A_{(1)} = \operatorname{diag}(1,0)$,
$A_{(2)} = \begin{psmallmatrix} 0 & 1\\ 2 & 0\end{psmallmatrix}$ and
$P = \operatorname{diag}(1,0)$ one finds
$\bar P A_{(2)} P = \begin{psmallmatrix} 0&0\\ 2&0\end{psmallmatrix}
\neq 0$, and the transformation creates a double pole at $v_2$.  With
$P$ as above the balance sends $\lambda \mapsto \lambda - 1$ and
$\mu \mapsto \mu + 1$.  It leaves the exponents at every other
singular point unchanged --- the residue at $v_k$ is conjugated by the
constant matrix $T(v_k)$ --- and the residue at infinity untouched,
since $T(\infty) = 1$.  A balance necessarily acts at two points and
never at one: for a Fuchsian connection
$A = \sum_k A_k\,\dd v/(v-v_k)$ the residue at infinity is
$-\sum_k A_k$, so the sum of all residues, that point included,
vanishes identically and no transformation can change an exponent at
a single point.

The target of the sequence is that every residue eigenvalue, which
for these systems has the form $n + a\eps$ with $n \in \Z$, be
normalized to $a\eps$: the integer parts must all be removed.  A
balance is admissible when $\lambda$ has integer part $n_1 > 0$ at
$v_1$ and $\mu$ has integer part $n_2 < 0$ at $v_2$, and eigenvectors
with $\eta^{\mathsf T}u \neq 0$ exist; it then decreases the
non-negative integer $\sum \lvert n \rvert$, summed over all residue
eigenvalues at all singular points including infinity, by two.  The
sequence therefore terminates.  What can fail is the existence of an
admissible pair at some step --- the partner with an integer part of
the opposite sign is guaranteed only in total, by the vanishing of
the sum of all residues, not at a point where the eigenvector
condition also holds.

\begin{example}[A balance and the constant step]
\label{ex:balance}
Take the one-variable Fuchsian system with singular points $0$ and
$1$,
\begin{equation}
  A \;=\; \frac{\dd v}{v}
  \begin{pmatrix} \eps & 0\\ 0 & 1+\eps\end{pmatrix}
  \;+\;
  \frac{\dd v}{v-1}
  \begin{pmatrix} \eps & 0\\ c & -1+\eps\end{pmatrix} .
\end{equation}
The exponent $1+\eps$ at $v=0$ is too high by one and the exponent
$-1+\eps$ at $v=1$ too low by one: an admissible pair.  A left
eigenvector of the first residue for $1+\eps$ is
$\eta^{\mathsf T} = (0,1)$; a right eigenvector of the second for
$-1+\eps$ is $u = (0,1)^{\mathsf T}$, since
$\big(A_{(2)}-(-1+\eps)\big)u = \begin{psmallmatrix}1&0\\c&0
\end{psmallmatrix}u = 0$; and $\eta^{\mathsf T}u = 1 \neq 0$.  Hence
$P = \operatorname{diag}(0,1)$, both conditions
\eqref{eq:balanceconditions} hold because the products are diagonal
or vanish, and \eqref{eq:balance} is
$T = \operatorname{diag}\big(1,\, v/(v-1)\big)$.  A diagonal $T$ acts
by $A'_{ij} = A_{ij}\,t_j/t_i$ off the diagonal and subtracts
$\dd\log t_i$ on it, so
\begin{equation}
  A' \;=\; \frac{\dd v}{v}
  \begin{pmatrix} \eps & 0\\ c & \eps\end{pmatrix}
  \;+\;
  \frac{\dd v}{v-1}
  \begin{pmatrix} \eps & 0\\ 0 & \eps\end{pmatrix} ,
\end{equation}
using $\frac{1+\eps}{v} + \frac{-1+\eps}{v-1}
- \big(\frac1v - \frac1{v-1}\big) = \frac{\eps}{v} +
\frac{\eps}{v-1}$ for the second diagonal entry and
$\frac{c}{v-1}\cdot\frac{v-1}{v} = \frac{c}{v}$ for the lower left.
All exponents are now $\eps$, but the entry $c$ is not proportional
to it.  The constant transformation
$T_c = \operatorname{diag}(\eps, 1)$ completes the reduction:
\begin{equation}
  A'' \;=\; T_c^{-1}A'T_c \;=\; \eps\left[
  \begin{pmatrix} 1 & 0\\ c & 1\end{pmatrix}\dd\log v
  \;+\;
  \begin{pmatrix} 1 & 0\\ 0 & 1\end{pmatrix}\dd\log(1-v)
  \right] ,
\end{equation}
the canonical form, with constant residues and the alphabet
$\{v,\,1-v\}$.
\end{example}

\begin{example}[Irregular and removable higher poles]
\label{ex:fuchs}
Consider first the scalar equation
\begin{equation}
  \frac{\dd F}{\dd w} = \left(\frac{\alpha}{w^2} +
  \frac{\beta\,\eps}{w}\right) F .
\end{equation}
The substitution $F = e^{-\alpha/w}\, G$ does remove the double
pole, since $\frac{\dd}{\dd w} e^{-\alpha/w} =
\frac{\alpha}{w^2}\,e^{-\alpha/w}$, leaving
$G' = (\beta\eps/w)\,G$ and $G = C\,w^{\beta\eps}$.  But
$e^{-\alpha/w}$ is not a rational --- not even a meromorphic ---
change of basis, and no rational one can do the job: a rational
$F = r(w)\,G$ shifts the connection by $-\dd\log r$, which has only
simple poles.  A scalar double pole is therefore a genuinely
irregular singularity, and no canonical form exists.

What \emph{is} removable is a higher pole sitting in a nilpotent
direction of a matrix residue.  For
\begin{equation}
  A \;=\; \frac{1}{w^2}
  \begin{pmatrix} 0 & \alpha \\ 0 & 0 \end{pmatrix} \dd w ,
  \qquad
  T = \begin{pmatrix} 1 & 0 \\ 0 & w \end{pmatrix}
  \quad\Longrightarrow\quad
  A' \;=\; \frac{1}{w}
  \begin{pmatrix} 0 & \alpha \\ 0 & -1 \end{pmatrix} \dd w ,
\end{equation}
a simple pole, reached by a rational transformation --- Fuchsian,
and ready for the balances above.  In the present calculation the
genuinely irregular case never occurs once the variables of
Section~\ref{sec:roots} are chosen correctly.
\end{example}

\subsection{The second variable}
\label{sec:secondvariable}

Moser reduction and balances are constructions for an ordinary
differential equation in one variable; the systems here have two.
The way through is to treat the variables in turn, and it is the
second step, not the first, that decides whether a canonical form is
reached.

Fix $w$ as a symbol and read the $v$-component of the Pfaffian
system,
\begin{equation}
  \partial_v \vec I \;=\; A_v(v,w;\eps)\, \vec I ,
\end{equation}
as an ordinary differential equation in $v$ over the field
$\Q(w,\eps)$ --- or over a finite algebraic extension of it, since
the singular points of this equation are the $v$-roots of the
letters, which are algebraic functions of $w$: for the letter
$1-v-w$ the root is $v = 1-w$, for $1-4vw$ it is $v = 1/(4w)$, and
for a letter quadratic in $v$ they are a conjugate pair.  Everything
of Section~\ref{sec:diagonal} applies verbatim over that field, the
spectator $w$ being carried along symbolically, and produces a
transformation $T_1(v,w;\eps)$ with
\begin{equation}
  T_1^{-1}A_v T_1 - T_1^{-1}\partial_v T_1
  \;=\; \eps \sum_a R_a\,\partial_v \log\varphi_a ,
  \qquad \partial_v R_a = 0 .
\end{equation}
The $R_a$ here are independent of $v$ --- elements of $\Q(w,\eps)$,
not yet constants in the sense of the canonical-form definition; the
residual dependence on the regulator is removed at the end, in
Section~\ref{sec:constT}.  Nothing has been required of the
$w$-direction, and in general
\begin{equation}
  A_w' \;=\; T_1^{-1}A_w T_1 - T_1^{-1}\partial_w T_1
\end{equation}
is not canonical.

The freedom left is a $v$-independent gauge $U(w;\eps)$.  Since
$\partial_v U = 0$, replacing $T_1$ by $T_1 U$ leaves the
$v$-direction free of an inhomogeneous term and merely conjugates its
residues,
\begin{equation}
  A_v'' \;=\; \eps \sum_a \big(U^{-1}R_a U\big)\,
  \partial_v\log\varphi_a ,
  \qquad
  A_w'' \;=\; U^{-1}A_w'\,U \;-\; U^{-1}\partial_w U ,
\end{equation}
while in the $w$-direction it acts in full.  In the canonical form
one and the same residue matrix stands in both directions, so two
requirements fall on $U$.  The conjugated residues must again be
constant,
\begin{equation}
  U^{-1} R_a\, U \;=\; R_a'' , \qquad
  R_a'' \ \text{independent of } w ,
  \label{eq:secondvaralg}
\end{equation}
an algebraic condition; and the $w$-connection must be canonical with
those same residues, which, written without the inverse, is
\begin{equation}
  \partial_w U
  \;=\; A_w'\,U \;-\; \eps\, U \sum_a R_a''\,
        \partial_w\log\varphi_a ,
  \label{eq:secondvar}
\end{equation}
a \emph{linear} equation for $U$ once the $R_a''$ are fixed.  The pair
\eqref{eq:secondvaralg}, \eqref{eq:secondvar} is bilinear in
$(U, R_a'')$, as the ansatz of Section~\ref{sec:diagonal} and the
conjugation of Section~\ref{sec:constT} are in their unknowns, and is
resolved the same way: once the $R_a''$ are fixed from a finite list
of candidates, both conditions are linear in $U$ and the pair is one
nullspace computation over $\Q(w,\eps)$, of finite size, carried out
once.  Flatness of the original system is a necessary condition for a
solution to exist and is not a sufficient one:
whether the pair is solvable is decided by the linear solve itself,
and nothing here asserts that the $w$-direction follows
automatically from the $v$-direction.

Two features of the outcome are worth recording.  The balances of
Section~\ref{sec:diagonal} introduce apparent singularities in $v$ at
$\eps$-dependent points; in the $w$-direction these appear as pure
scalar $\dd\log$ terms with integer residues, and a rational scalar
gauge removes them, so that they do not survive in the final
alphabet.  And the acceptance criterion is the two-variable identity:
the composed transformation is admitted only when
\eqref{eq:epsform} holds in \emph{both} directions with one and the
same set of constant residues.  A one-variable check establishes
nothing about a two-variable system.

\subsection{The entries below the diagonal: the completion equation}
\label{sec:completion}

Suppose the diagonal blocks are already canonical --- each equal to
$\eps\,e_k$ with $e_k$ a constant-residue logarithmic form ---
while the entries below the diagonal, denoted $b_{kj}$ with $j<k$,
are not.  They are brought to canonical form one block row at a
time, by transformations of unipotent shape:
\begin{equation}
  T = 1 + D, \qquad
  D\ \text{supported on block row } k,\ \text{columns } j<k .
\end{equation}
Because $D$ is strictly triangular, $D^2 = 0$ and $T^{-1} = 1 - D$;
moreover $D A D$ and $D\,\dd D$ vanish --- the latter because $D$
and $\dd D$ share the same support, the former because $A$ is block
lower triangular, so its $(j,k)$ entries with $j < k$ are zero ---
and the transformation law \eqref{eq:gauge} is exact without
expansion:
\begin{equation}
  A' \;=\; A + A D - D A - \dd D .
\end{equation}
Reading off the $(k,j)$ block of the requirement that $A'$ be
canonical gives the \emph{completion equation}
\begin{equation}
  \dd D_{kj}
  \;=\; \eps\big( e_k\, D_{kj} - D_{kj}\, c_j \big)
  \;+\; \bar b_{kj}
  \;-\; \eps \sum_a K_a\, \dd\log\varphi_a ,
  \label{eq:completion}
\end{equation}
in which $c_j$ is the canonical form of the diagonal block in
column $j$ (so that $a_{jj} = \eps\,c_j$, in parallel with
$a_{kk} = \eps\,e_k$), the source
$\bar b_{kj} = b_{kj} - \sum_{j<m<k} D_{km}\, b_{mj}$ collects the
original coupling together with the contributions of blocks of the
same row completed previously, and the matrices $K_a$ --- constant in
the kinematics, though at this point still allowed to depend on the
regulator, which Section~\ref{sec:constT} removes --- are unknowns
alongside $D_{kj}$: they become the residues of the completed entry.
The essential feature of \eqref{eq:completion} is its structure in
the unknowns $(D_{kj}, K_a)$: it is \emph{affine-linear} --- a linear
operator applied to the unknowns, plus the fixed inhomogeneous source
$\bar b_{kj}$.  The standard consequence follows at once: its
solution set is an affine subspace, one particular solution plus the
null space of the linear part.  That null space is exactly the
non-uniqueness of Remark~\ref{rem:representative}, and the same
affine structure will appear once more, in the constraint system of
Section~\ref{sec:kernel}.

\begin{example}[A completion in one variable]
\label{ex:completion}
In one variable, with diagonal data $e = \dd\log v$ and $c = 0$,
take the coupling
\begin{equation}
  b \;=\; \eps\,\dd v \;+\; \frac{\dd v}{1+v} .
\end{equation}
The two pieces are absorbed by different mechanisms.  For the first,
try a rational $D$: with $D = \dfrac{\eps\, v}{1-\eps}$ one has
\begin{equation}
  \eps\, e\, D + \eps\,\dd v
  = \eps\,\frac{\dd v}{v}\cdot\frac{\eps v}{1-\eps} + \eps\,\dd v
  = \frac{\eps}{1-\eps}\,\dd v
  = \dd D ,
\end{equation}
so the non-logarithmic part of the coupling disappears into the
transformation.  The second piece is itself a logarithmic form with
constant coefficient, $\dd\log(1+v)$: no rational $D$ can absorb it
(that would require $\log(1+v)$ itself), and none should --- it is
kept as a residue.  Matching it in \eqref{eq:completion} requires
$\eps K_{1+v} = 1$, so the completed entry carries
$\dd\log(1+v)$ at order $\eps^0$ instead of at order $\eps$.  That
mismatch is a normalization of the basis, constant in the kinematics,
and is removed by the final transformation of
Section~\ref{sec:constT}.  Rational parts of a coupling are absorbed
by $D$; logarithmic parts with constant coefficients become
residues; normalization mismatches in $\eps$ are deferred to the
end --- the same three cases exhaust the general situation.
\end{example}

\subsection{Solution of the completion equation by modular
arithmetic}
\label{sec:modular}

For the systems of interest the unknown $D_{kj}$ may have several
hundred rational-function entries of high degree, and solving
\eqref{eq:completion} symbolically is prohibitive.  Being affine-linear,
it
can instead be solved over finite fields and reassembled exactly.
One writes
$D_{kj} = N(v,w;\eps)\big/\prod_a \varphi_a^{m_a}$, with the
denominator a product of letters and the numerator a polynomial with
undetermined coefficients.  The ansatz is bounded by the source.
Suppose $\bar b_{kj}$ has a pole of order $p_a$ at $\varphi_a = 0$
while $D_{kj}$ has one of order $m_a > p_a$, with leading coefficient
$D^{(m)}$.  At order $m_a + 1$ in that pole only two terms of
\eqref{eq:completion} contribute, $\dd D_{kj}$ and
$\eps(e_k D_{kj} - D_{kj}c_j)$, and they must cancel:
\begin{equation}
  -\,m_a\, D^{(m)} \;=\; \eps\big(R^e\,D^{(m)} - D^{(m)}\,R^c\big) ,
\end{equation}
with $R^e, R^c$ the residues of the diagonal data at that letter.
Thus $D^{(m)}$ is an eigenvector of the linear map
$X \mapsto \eps\,(R^eX - XR^c)$ with eigenvalue $-m_a$; the
eigenvalues of that map are $\eps(\rho_i - \rho'_j)$ with
$\rho_i, \rho'_j$ the eigenvalues of $R^e$ and $R^c$, and for generic
$\eps$ none of them is a nonzero integer.  Hence $D^{(m)} = 0$ and
\begin{equation}
  m_a \;\le\; p_a ,
  \qquad
  \deg N \;\le\; \sum_a m_a \deg\varphi_a
  \;+\; \deg\bar b_{kj} \;+\; 1 ,
\end{equation}
the second a degree count in the same spirit, with
$\deg \bar b_{kj}$ the excess of the numerator degree of the source
over its denominator degree.  The degree bound is a bound and not a
value: if the fit is inconsistent it is raised and the fit repeated,
and inconsistency is exactly what an over-determined linear system
reports.  Substituting rational values of $(v,w)$ \emph{and of
the regulator}, all reduced modulo a large prime $p$, turns
\eqref{eq:completion} into exact linear equations over the finite
field $\Fp$ for the unknown coefficients; enough substitutions
determine them modulo $p$.  The dependence on $\eps$ is recovered by
the same device: the solution is a rational function of the
regulator, and a rational function with numerator and denominator
degrees $d_N$ and $d_D$ is determined by $d_N + d_D + 1$ samples ---
a polynomial of degree $d$ by $d+1$ --- so one samples until the
reconstruction stabilizes and withholds a further handful of
regulator values as a check.  Repeating over several primes,
combining by the Chinese remainder theorem, and lifting each
coefficient to a rational number by rational reconstruction yields
an exact candidate, confirmed on evaluations withheld from the fit
and at a fresh prime.  None of this is a proof; the proof is the
substitution of the reconstructed $(D, K)$ back into
\eqref{eq:completion} --- or into the equivalent statement for the
fully assembled family --- and the exact verification of the
identity.

\begin{example}[Rational reconstruction]
Take $p = 101$ and suppose a coefficient is determined as
$29 \bmod 101$.  Rational reconstruction seeks a fraction of small
numerator and denominator congruent to it: since
$7 \cdot 29 = 203 = 2\cdot 101 + 1 \equiv 1 \pmod{101}$, one has
$29 \equiv \tfrac{1}{7} \pmod{101}$.  A second prime, $p = 103$,
where $\tfrac17 \equiv 59 \pmod{103}$ (as $7 \cdot 59 = 413 =
4\cdot103+1$), and the Chinese remainder lift of $(29, 59)$ modulo
$101\cdot103$ reconstructs $\tfrac17$ again.  The guarantee behind
the search is a lattice bound: modulo a single prime $p$ the
reconstruction is unique only for numerator and denominator below
$\sqrt{p/2}$ --- about $7.1$ for $p = 101$, so $\tfrac17$ is
admissible there only just, which is exactly why the second prime is
not optional.
\end{example}

\begin{remark}[Non-uniqueness of the completion]
\label{rem:representative}
When the linear part of \eqref{eq:completion} has a null space the
completed entry is not unique: the affine solution set noted above is
then a positive-dimensional affine subspace.  Two solutions differ by
an element of that null space, and the canonical forms they produce
are related by a constant
transformation that preserves the diagonal blocks: the two are
equally canonical, but their residues $K_a$ are not the same
matrices.  What the choice does affect materially is the degrees of
the rational functions inserted into the sources $\bar b$ of the rows
still to be completed, and degree growth compounds from row to row.
A definite and favourable choice results from requiring the
coefficients of the highest admissible numerator monomials to vanish,
which selects a representative of low degree.
\end{remark}

\subsection{The residual dependence on the regulator}
\label{sec:constT}

After the completion, the connection is logarithmic with
$\eps$-independent letters, but its residue matrices may still
depend on $\eps$ --- rational functions of $\eps$ in place of the
required constants times $\eps$.  The origin of this residue is the
normalization of the master basis: each sector's masters were
defined with $\eps$-dependent prefactors, and a change of
normalization is a transformation \emph{constant in the
kinematics}.  For constant $T(\eps)$ the inhomogeneous term of
\eqref{eq:gauge} vanishes and the transformation acts by
simultaneous conjugation of all residues:
\begin{equation}
  T^{-1}(\eps)\, \tilde R_a(\eps)\, T(\eps) \;=\; \eps\, R_a
  \qquad \text{for all } a .
\end{equation}
Written without the inverse, $\tilde R_a(\eps)\,T(\eps) = \eps\,
T(\eps)\,R_a$, the condition is linear in $T$ once the target
residues are fixed --- a finite-dimensional problem, solved at
sampled rational values of the regulator and interpolated, with the
cheap necessary test that the eigenvalues of
$\tilde R_a(\eps)/\eps$ be $\eps$-independent applied first.  When a
solution exists it is unique up to right multiplication by a matrix
commuting with every $R_a$.  Composing the class transformations, the
row completions, and this constant conjugation yields the matrix $T$
carrying the derived system \eqref{eq:desys} into the canonical form
\eqref{eq:epsform}.

%======================================================================
\section{Algebraic letters: square roots and their rationalization}
\label{sec:roots}

\subsection{Square roots from the geometry of phase space}

Nothing so far required the letters to be polynomials, and for a
minority of the blocks they are not.  The origin is the geometry of
the momenta.  Eliminating the direction of an unobserved momentum at
fixed invariants requires solving quadratic equations, and the
boundary of the region in which such an equation has real solutions
is the vanishing of a Gram determinant.  For a sub-configuration
carrying three momenta $P_1$, $P_2$, $P_3 = -P_1-P_2$ that Gram
determinant is a K\"all\'en polynomial: with
$2P_1\!\cdot\!P_2 = P_3^2 - P_1^2 - P_2^2$,
\begin{equation}
  \det\begin{pmatrix}
    2P_1^2 & 2P_1\!\cdot\!P_2\\
    2P_1\!\cdot\!P_2 & 2P_2^2
  \end{pmatrix}
  = 4P_1^2P_2^2 - \big(P_3^2-P_1^2-P_2^2\big)^2
  = -\,\lambda\big(P_1^2, P_2^2, P_3^2\big) ,
  \label{eq:gramkallen}
\end{equation}
with $\lambda(a,b,c) = a^2+b^2+c^2-2ab-2bc-2ca$ the triangle
function, for which $\lambda(a,b,c) = (c-a-b)^2 - 4ab$.  The
sub-configurations of the double-real channel that produce roots
carry the virtualities $s$, $-t$, $-u$, and $\lambda$ is homogeneous
of degree two, so that
$\lambda(s,-t,-u)/s^2 = \lambda(1,v,w)$: dividing
\eqref{eq:gramkallen} by $s^2$ gives the polynomial that governs
them,
\begin{equation}
  \lambda_1(v,w) \;=\; \lambda(1, v, w)
  \;=\; (1 - v - w)^2 - 4\,v\,w .
  \label{eq:kallen}
\end{equation}
Its zero locus lies \emph{inside} the physical triangle
\eqref{eq:region} rather than on its boundary.  Indeed
\begin{equation}
  \lambda_1 \;=\;
  \Big(1 - \big(\sqrt v + \sqrt w\big)^2\Big)\,
  \Big(1 - \big(\sqrt v - \sqrt w\big)^2\Big) ,
  \label{eq:kallenfactored}
\end{equation}
and the second factor is positive throughout the triangle, so
$\lambda_1$ vanishes exactly on $\sqrt v + \sqrt w = 1$ --- through
the interior point $v = w = \tfrac14$ --- and is positive on the side
of it containing the origin, where $\sqrt{\lambda_1}$ is real.  In
the energy and angle of \eqref{eq:vwangle}, with
$\xi = v+w = 2E_3/\sqrt s$ and $4vw = \xi^2\sin^2\theta$, the locus is
$\xi\,(1+\lvert\sin\theta\rvert) = 1$ --- an unremarkable
configuration of the final state, at which the masters of the sectors
sensitive to it nevertheless acquire a branch point of square-root
type: a pseudo-threshold.  A square-root branch point cannot be
produced by logarithmic differentials of polynomials alone: the
alphabet must contain $\sqrt{\lambda_1}$, typically through letters
of the form
\begin{equation}
  \frac{1 - v - w - \sqrt{\lambda_1}}
       {\,1 - v - w + \sqrt{\lambda_1}\,} ,
\end{equation}
which are exchanged with their reciprocals under the sign change
$\sqrt{\lambda_1} \to -\sqrt{\lambda_1}$ (in the language of
covering spaces, the deck transformation of the double cover).  The
double-real channel realizes six such radicands:
\begin{equation}
  \lambda_1,\qquad
  \lambda_2 = \lambda_1(-v,w),\qquad
  \lambda_3 = \lambda_1(v,-w),\qquad
  4v + w^2,\qquad v^2 + 4w,\qquad 1 - 4vw .
  \label{eq:radicands}
\end{equation}
Under the involution $v \leftrightarrow w$ --- the exchange of the
incoming partons, used in Definition~\ref{def:class} --- the pairs
$\{\lambda_2, \lambda_3\}$ and $\{4v{+}w^2,\, v^2{+}4w\}$ are
exchanged, while $\lambda_1$ and $1-4vw$ are invariant; in the chart
of Example~\ref{ex:kallenchart} the involution reads
$(x,y)\mapsto(1-x,1-y)$.

\begin{definition}[Root content of a system]
Radicands are counted modulo perfect-square factors and modulo
products: $\sqrt{q_1}$, $\sqrt{q_2}$ and $\sqrt{q_1 q_2}$ constitute
\emph{two} independent roots, since the third is determined by the
other two.  Formally, the radicands of a system generate a subgroup
of the square-class group
$\Q(v,w)^\times / \big(\Q(v,w)^\times\big)^2$, and the \emph{root
content} is the rank of that subgroup.
\end{definition}

The rank is computed by elementary linear algebra over
$\mathbb{F}_2$.  Factor each radicand into irreducible polynomials,
$q_i = \prod_\alpha \pi_\alpha^{\,n_{i\alpha}}$, and record only the
parities: a radicand is a square exactly when every $n_{i\alpha}$ is
even, and a product of radicands corresponds to the sum of their
exponent vectors.  The root content is therefore the rank over
$\mathbb{F}_2$ of the matrix
\begin{equation}
  \big(n_{i\alpha} \bmod 2\big)_{i\alpha} ,
\end{equation}
one row per radicand, one column per irreducible factor.  For
$\lambda_1$, $\lambda_2$, $\lambda_1\lambda_2$ the third row is the
sum of the first two, and the rank is two.

The radicands themselves are read off the derived connection, before
any transformation of it: the discriminants of the irreducible
quadratic factors of its singular loci, together with those occurring
in the eigenvalues of its residue matrices.  What this yields is an
\emph{upper bound} on the root content, not its value: a quadratic
factor may be an apparent singularity, removable by a transformation,
and its discriminant then contributes no root to the final alphabet.
The bound is what selects a first candidate among the constructions
below; a smaller root content is discovered when the construction
succeeds without the corresponding root.  In the double-real channel
about half of the families involve no root at all; a third involve
one; a dozen involve two; and three families --- the hard case
treated in Section~\ref{sec:galois} --- involve three.

\subsection{Rationalization by change of variables}

\begin{definition}[Rationalizing substitution]
A \emph{rationalizing substitution} for a radicand $q(v,w)$ is a
rational change of variables $(v,w) = \big(f(x,y),\,g(x,y)\big)$
whose Jacobian determinant is not identically zero and for which
$q(f,g) = r(x,y)^2$ with $r$ a
rational function.  Under such a substitution $\sqrt q$ becomes
the rational function $r$, and a system whose alphabet involves only
$\sqrt q$ acquires polynomial letters in $(x,y)$.
\end{definition}

\begin{example}[Rationalizing the K\"all\'en root]
\label{ex:kallenchart}
Consider the substitution suggested by the factorized form of the
degenerate configurations,
\begin{equation}
  v = x\,y, \qquad w = (1-x)(1-y) .
\end{equation}
Then
\begin{equation}
  1 - v - w = x + y - 2xy,
  \qquad
  \lambda_1 = (x+y-2xy)^2 - 4\,xy\,(1{-}x)(1{-}y)
  = (x-y)^2 ,
\end{equation}
by direct expansion.  Hence $\sqrt{\lambda_1} = x - y$, the
Jacobian determinant is $\partial(v,w)/\partial(x,y) = x - y$,
non-vanishing away from the diagonal, and the algebraic letter above
becomes
\begin{equation}
  \frac{(x+y-2xy)-(x-y)}{(x+y-2xy)+(x-y)}
  \;=\; \frac{y\,(1-x)}{x\,(1-y)} ,
\end{equation}
a ratio of polynomial letters.  Two facts about the substitution
deserve recording.  It is two-to-one: both $v$ and $w$ are symmetric
under $x \leftrightarrow y$, and that exchange is precisely the deck
transformation $\sqrt{\lambda_1} \to -\sqrt{\lambda_1}$, branched on
the diagonal $x = y$ where the Jacobian vanishes.  And its real
image is not the whole triangle \eqref{eq:region} but the
sub-region $\sqrt v + \sqrt w \le 1$, i.e.\ $\lambda_1 \ge 0$;
outside it $\sqrt{\lambda_1}$ is imaginary, and reaching it is an
analytic continuation to be performed deliberately.  Paths chosen in
$(x,y)$ therefore live in that sub-region.  The sign variants
$v = \pm xy$, $w = \mp(1-x)(1-y)$ rationalize $\lambda_2$ and
$\lambda_3$ in the same way, and substitutions of the same conic
type rationalize $4v+w^2$, $v^2+4w$ and $1-4vw$ --- one
substitution for each of the six radicands.
\end{example}

The construction behind all such substitutions is classical: the
equation $q(v,w) = r^2$ defines a conic in suitable projective
coordinates, and a conic possessing one rational point is rationally
parametrized by the pencil of lines through that point --- the same
device that parametrizes the Pythagorean triples.  In the case needed
here the rational point can always be taken at infinity.  Regard $q$
as a quadratic in one of the variables,
\begin{equation}
  q \;=\; a(v)\,w^2 + b(v)\,w + c(v) ,
  \qquad a = \alpha^2 \ \text{a perfect square} ,
\end{equation}
and set $r = \alpha w + \tau$ with $\tau$ a new variable.  The
squares $\alpha^2w^2$ cancel from $r^2 = q$, leaving a
\emph{linear} equation for $w$:
\begin{equation}
  2\alpha\,w\,\tau + \tau^2 = b\,w + c
  \qquad\Longrightarrow\qquad
  w = \frac{c - \tau^2}{2\alpha\tau - b},
  \qquad
  \sqrt q = \alpha\,w + \tau .
  \label{eq:pencil}
\end{equation}
The pair $(v,\tau)$ is a rationalizing chart for $q$, and the
identity $q = r^2$ in it is a matter of expansion.  If the leading
coefficient is not a square, one exchanges the roles of $v$ and $w$,
or shifts $w$ so that the constant term $c$ becomes a square and
applies \eqref{eq:pencil} to the reciprocal variable.  The systematic
treatment of such substitutions is that of Besier, van Straten,
Wasser and Weinzierl~\cite{besier}.
The symmetric chart of Example~\ref{ex:kallenchart} is an equivalent
choice adapted to the involution $v \leftrightarrow w$: for
$\lambda_1$ as a quadratic in $w$ one has $\alpha = 1$,
$b = -2(1+v)$, $c = (1-v)^2$, and \eqref{eq:pencil} rationalizes it
just as well.

Working in the new variables requires one transformation of the
whole system by the chain rule,
\begin{equation}
  A_x = A_v\,\partial_x f + A_w\,\partial_x g,
  \qquad
  A_y = A_v\,\partial_y f + A_w\,\partial_y g ,
\end{equation}
with flatness re-established in $(x,y)$; afterwards the
constructions of Sections~\ref{sec:epsform}--\ref{sec:finding}
apply verbatim.

\subsection{Several roots}

A system whose alphabet involves two independent roots requires a
substitution rationalizing both at once, and such substitutions are
found by iterating \eqref{eq:pencil}.  Rationalize the first root;
under that substitution the second radicand becomes a polynomial
$q_2(x,y)$ in the new variables, and \eqref{eq:pencil} applied to
$q_2$ in one of them rationalizes it as well, leaving the first root
--- already rational --- untouched.  Two conditions must hold for the
step to succeed, and neither is automatic.  The pullback of a conic
under a map of degree two is in general a quartic, so $q_2$ must
still be of degree two in the variable chosen; and its leading
coefficient in that variable must be a perfect square, so that the
rational point at infinity used by \eqref{eq:pencil} is available.
The choice of variable is therefore not free, and for a general pair
of radicands the construction stops here.  For the K\"all\'en pairs
of the double-real channel both conditions are met.

\begin{example}[Two K\"all\'en roots at once]
\label{ex:tworoots}
Start from the chart of Example~\ref{ex:kallenchart}, in which
$\sqrt{\lambda_1} = x-y$.  In the same chart the second radicand
$\lambda_2 = \lambda_1(-v,w) = (1+v-w)^2 + 4vw$ becomes, using
$1 + v - w = x+y$,
\begin{equation}
  \lambda_2 \;=\; (x+y)^2 + 4\,xy\,(1-x)(1-y)
  \;=\; (1-2x)^2\,y^2 \;+\; 2x\,(3-2x)\,y \;+\; x^2 ,
\end{equation}
a quadratic in $y$ whose leading coefficient is the square
$(1-2x)^2$.  Its discriminant, $32\,x^2(1-x)$, is not a square, so
$\lambda_2$ does not factor into rational linear pieces; but that is
not what is needed.  Applying \eqref{eq:pencil} with $\alpha = 1-2x$,
$b = 2x(3-2x)$, $c = x^2$ and the new variable $\tau$,
\begin{equation}
  y \;=\; \frac{x^2 - \tau^2}
  {2\big[(1-2x)\,\tau - x\,(3-2x)\big]} ,
  \qquad
  \sqrt{\lambda_2} \;=\; (1-2x)\,y + \tau ,
\end{equation}
and in the chart $(x,\tau)$ both roots are rational: $v = xy$,
$w = (1-x)(1-y)$ and $\sqrt{\lambda_1} = x - y$ with $y$ the rational
function just displayed.  The check is quick at $x = \tfrac12$, where
$\lambda_2 = 2y + \tfrac14$ and the formula gives
$y = \tfrac12(\tau^2 - \tfrac14)$, so that $\lambda_2 = \tau^2$
indeed.  Three substitutions of this kind --- one for each pair of
K\"all\'en-type radicands --- cover every two-root family that
occurs, and the canonical connections of those families are written
in the final pair of parameters.
\end{example}

\subsection{The Galois algebra of the roots}
\label{sec:galois}

For the families involving three independent roots, no rational
substitution rationalizes all three simultaneously.  This follows
from the geometry of the covering surface, not from a failure of
search, and the argument is worth carrying out, since it is the only
place in these notes where a construction is ruled out rather than
performed.

A substitution rationalizing $q_1, q_2, q_3$ at once is a rational
map from the plane onto the surface $X$ on which all three roots are
single-valued: the $(\Z/2)^3$ cover of the $(v,w)$-plane branched
over the three conics $q_i = 0$.  Adjoining the three roots as
coordinates gives $X$ a projective model in $\mathbb{P}^5$, in
homogeneous coordinates $(v,w,h,r_1,r_2,r_3)$,
\begin{equation}
  r_i^2 \;=\; Q_i(v,w,h) \;=\; h^2\,q_i(v/h,\, w/h) ,
  \qquad i = 1,2,3 ,
\end{equation}
with $Q_i$ the degree-two homogenization of $q_i$, so that each
equation is a genuine quadric.  The model is a complete intersection
of three quadrics, of degree $2^3 = 8$, matching the eight sheets of
the cover.  Its singularities sit over the points where two branch
conics meet.  The pairwise intersections are transverse except for
finitely many simple tangencies, and each tangency lifts to two
\emph{ordinary double points} of the cover --- points at which the
surface looks locally like $xy = z^2$, the simplest surface
singularity, called a node; in suitable local coordinates the lift is
the equation $r_2^2 - r_1^2 + z^2 = 0$.  How many there are is
immaterial, because nodes are rational double points and their
resolution is \emph{crepant}: the canonical class pulls back without
correction, $K_{\widetilde X} = K_X$.  For a complete intersection of
three quadrics in $\mathbb{P}^5$ the adjunction formula gives
$K = \mathcal{O}\big(\sum_i d_i - n - 1\big) =
\mathcal{O}(6 - 6)$, trivial.  A surface with trivial canonical
class, vanishing irregularity
$h^1(\widetilde X, \mathcal{O}_{\widetilde X}) = 0$ and geometric genus
$p_g = \dim H^0(\widetilde X, K_{\widetilde X}) = 1$ is by definition
a K3 surface, and $\widetilde X$ has these invariants.

The conclusion follows from two classical facts, for which
Beauville's book~\cite{beauville} is the reference.  A surface is
\emph{unirational} if some rational map from the plane onto it is
dominant --- which is exactly what a simultaneous rationalizing
substitution would provide.  By Castelnuovo's theorem, in
characteristic zero a unirational surface is rational; and a rational
surface has $p_g = 0$, since $p_g$ is a birational invariant and the
plane has no holomorphic two-forms.  But $\widetilde X$ has
$p_g = 1$.  No such substitution exists.

The response is to renounce parametrization and compute in the
extension the roots generate.  Adjoin them as symbols:
\begin{equation}
  E \;=\; \Q(v,w,\eps)\,[\,r_1,r_2,r_3\,]\ \big/\
  \big(\,r_i^2 - q_i\,\big) ,
\end{equation}
an algebra of dimension $2^3$ with basis
$r_S = \prod_{i\in S} r_i$ indexed by subsets
$S \subseteq \{1,2,3\}$; when the $q_i$ are independent square
classes --- which is what root content three means --- $E$ is a
field extension of degree eight.  Its Galois group is $(\Z/2)^3$,
acting by the independent sign changes $r_i \mapsto -r_i$.  Every
element decomposes into the characters of this group (its
\emph{grades}); multiplication respects the grading,
\begin{equation}
  r_S\, r_T \;=\; \Big(\prod_{i \in S\cap T} q_i\Big)\;
  r_{S \mathbin{\triangle} T} ,
\end{equation}
and differentiation acts grade-wise through
$\dd r_S = \tfrac12\, r_S \sum_{i\in S} \dd\log q_i$.  All linear
problems of Section~\ref{sec:finding} --- the completion equation,
its evaluation over finite fields, the reconstruction --- can be
posed on the $8$-component coefficient vectors, whose entries are
ordinary rational functions.  Evaluations modulo a prime are
performed at points where every $q_i$ is a nonzero quadratic
residue, so that all eight embeddings of the algebra into $\Fp$
exist and the grade decomposition becomes a discrete Fourier
transform over $(\Z/2)^3$; agreement of a solution across all eight
sign embeddings is a consistency requirement that a correct
computation must, and does, satisfy.  No better description of these
systems is presently known.  What the construction requires of a
system is one checkable property: that its alphabet involve three
independent square classes drawn from the radicands
\eqref{eq:radicands}, so that the grading is by $(\Z/2)^3$ and the
rules above apply verbatim.  The three families of the channel that
have root content three have that property; carrying the construction
through on them is part of the work that Section~\ref{sec:proof}
records as remaining.  It has so far been carried through on model
systems of every root rank.

\begin{example}[One root, in components]
\label{ex:galois1}
The mechanism is already complete for $k=1$.  Let
$E = \Q(v)[r]/(r^2-q)$ with $q = q(v)$ not a square, so that
$\{1, r\}$ is a basis and every element is $c_0 + c_1 r$ with
$c_0, c_1$ rational.  Multiplication in this basis is the matrix
\begin{equation}
  (c_0 + c_1 r)\,(z_0 + z_1 r)
  \;=\; (c_0z_0 + c_1z_1 q) + (c_0z_1 + c_1z_0)\,r ,
  \qquad
  \begin{pmatrix} z_0\\ z_1\end{pmatrix}
  \longmapsto
  \begin{pmatrix} c_0 & c_1 q\\ c_1 & c_0 \end{pmatrix}
  \begin{pmatrix} z_0\\ z_1\end{pmatrix} ,
\end{equation}
whose determinant $c_0^2 - c_1^2q$ is the norm; it vanishes only for
$c_0 = c_1 = 0$, which is why $E$ is a field.  Differentiation is
grade-preserving up to the factor $r$: from $r^2 = q$ one has
$2r\,\dd r = \dd q$, that is
\begin{equation}
  \dd r \;=\; \tfrac12\, r\, \dd\log q .
\end{equation}
Consequently a one-form equation in $E$, say
$\dd D = \beta\,D + \sigma$ with $D = D_0 + D_1 r$,
$\beta = \beta_0 + \beta_1 r$, $\sigma = \sigma_0+\sigma_1 r$,
separates into two coupled equations with rational coefficients,
\begin{align}
  \dd D_0 &= \beta_0 D_0 + q\,\beta_1 D_1 + \sigma_0 ,\\
  \dd D_1 + \tfrac12 D_1\,\dd\log q
          &= \beta_0 D_1 + \beta_1 D_0 + \sigma_1 ,
\end{align}
the even and odd components.  Nothing algebraic has been lost and
nothing irrational remains in the coefficients: the price of the root
is a doubling of the number of unknowns.  For $k$ roots the same
computation gives $2^k$ components, indexed by the subsets $S$, with
the multiplication and differentiation rules displayed above.
\end{example}

%======================================================================
\section{The general solution: paths and polylogarithms}
\label{sec:paths}

The canonical form reduces integration to the evaluation of the
iterated integrals in \eqref{eq:pexp} along paths in the physical
region.  For letters \emph{linear} in the integration parameter the
iterated integrals belong to a standard class:

\begin{definition}[Generalized polylogarithms]
For $a_1,\dots,a_k \in \mathbb{C}$ define recursively
\begin{equation}
  G(a_1,a_2,\dots,a_k;\,t) \;=\;
  \int_0^{t}\frac{\dd t'}{t' - a_1}\;
  G(a_2,\dots,a_k;\,t') ,
  \qquad G(;t) = 1 ,
  \label{eq:gpl}
\end{equation}
the integration running along a path from $0$ to $t$ that avoids the
points $a_i$, with a branch of each logarithm fixed once along it; the
value depends on the homotopy class of that path and on nothing else.
The recursion is well posed unless every $a_i$ vanishes, in which case
the integral diverges at the lower limit; those words are defined
instead by $G(\vec 0_k; t) = \tfrac{1}{k!}\log^k t$, and the two
prescriptions cover every word.  These \emph{generalized
polylogarithms}~\cite{goncharov} of weight $k$ contain the classical
functions: $G(0,1;t) = -\operatorname{Li}_2(t)$, and all functions of
weight two and three reduce to classical polylogarithms.
\end{definition}

\begin{example}[From the definition to the dilogarithm]
Two applications of the definition, for $0 < t < 1$:
\begin{equation}
  G(1;t) = \int_0^t \frac{\dd t'}{t'-1} = \log(1-t),
  \qquad
  G(0,1;t) = \int_0^t \frac{\dd t'}{t'}\,G(1;t')
           = \int_0^t \frac{\dd t'}{t'}\,\log(1-t')
           = -\operatorname{Li}_2(t) .
\end{equation}
This is Example~\ref{ex:li2} read backwards: the two iterated
integrals generated there by the canonical recursion are precisely
$F_2^{(1)} = G(1;w)$ and $F_2^{(2)} = -G(0,1;w)$.
\end{example}

Evaluating a family is then a matter of choosing a base point $p_0$
in the triangle \eqref{eq:region} and a path to the point of
interest, composed of segments on which every letter restricts to a
polynomial of degree one in the path parameter.  The criterion is
exactly what the definition of $G$ requires: if
$\varphi_a\big(\gamma(t)\big) = \kappa_a\,(t - \alpha_a)$ then
\begin{equation}
  \dd\log\varphi_a\big(\gamma(t)\big)
  \;=\; \frac{\dd t}{t - \alpha_a} ,
\end{equation}
the constant $\kappa_a$ dropping out, so that every iterated integral
of \eqref{eq:pexp} along the segment is one of the functions $G$,
with indices the roots $\alpha_a$ of the restricted letters.
When a letter restricts to a quadratic, subdividing the segment does
not help: the restriction stays quadratic on every piece.  The
remedy is either to choose the segment differently --- in practice
segments parallel to the coordinate axes, on which the bilinear
letters produced by the rationalizing substitutions become linear
again --- or to factor the quadratic and admit its two linear
factors, with algebraic constants, as letters in their own right.
Flatness guarantees that the result depends only on the endpoint and
on the homotopy class of the path; the zero locus of a letter is
crossed only as a deliberate analytic continuation, never in
passing.  Families whose canonical form lives in rationalizing
variables are evaluated in those variables throughout, along paths
chosen by the same linearity criterion, and only the final values
are mapped back to $(v,w)$.

\begin{example}[Choosing the path]
\label{ex:path}
Suppose the alphabet of a family is $\{v,\ w,\ 1-v-w,\ 1-4vw\}$ and
the value at a point $(v_1,w_1)$ of the triangle is wanted.  Here
$1-4vw$ occurs as a polynomial letter in its own right; that the same
polynomial occurs elsewhere as a radicand, in \eqref{eq:radicands},
is a separate matter and plays no role in this example.  Take
the path in two segments parallel to the axes: first in $v$ at
fixed $w$, then in $w$ at fixed $v = v_1$.  On the second segment,
parametrized by $w = t$,
\begin{equation}
  v \to v_1, \qquad
  w \to t, \qquad
  1-v-w \to (1-v_1) - t, \qquad
  1-4vw \to 1 - 4 v_1 t :
\end{equation}
the letter $v$ is constant and drops out of the connection, and
every remaining letter has degree one in $t$, so the iterated
integrals along the segment are the functions
$G(a_1,\dots,a_k;\,t)$ with indices drawn from the letters' roots,
$a_i \in \{0,\ 1-v_1,\ 1/(4v_1)\}$, seeded by the values
accumulated on the first segment.  Had the diagonal segment
$v = w = t$ been chosen instead, the last letter would restrict to
$1-4t^2$ --- quadratic, and still quadratic on every sub-segment
--- and one would either return to axis-parallel segments or admit
its factors $1 \mp 2t$, in general with algebraic roots, as letters
in their own right.
\end{example}

\subsection{The order of expansion required}

The cost of an iterated integral grows steeply with its weight, so
one does not expand every master to the same depth; the required
depth is dictated by where the master enters the cross section.  The
rule is arithmetic.  The assembled combination is
\begin{equation}
  \sigma \;=\; \sum_m c_m(\eps)\, I_m(v,w;\eps),
  \qquad
  c_m(\eps) = \sum_{j \ge -p_m} \eps^{\,j} c_m^{(j)} ,
\end{equation}
with $p_m$ the order of the pole that the coefficient of the $m$-th
master carries after reduction and after the expansion of the
measure.  If $\sigma$ is wanted through order $\eps^{\,n}$, then
$I_m$ is wanted through order $\eps^{\,n + p_m}$: the deepest pole a
master multiplies anywhere in the combination sets its depth.  Thus
$\eps^{-2}\big(w^{\eps} - 1\big) = \eps^{-1}\log w +
\tfrac12\log^2 w + \tfrac16\,\eps\log^3 w + \dots$ needs the
$\eps^{2}$ coefficient of $w^\eps$ for its finite part and no more.
Cancellations of poles between terms do not relax the requirement,
since each master must be known before the cancellation can be
exhibited.  In the double-real
channel the outcome is lopsided: of the $347$ masters, $203$ are
needed only to their finite order and $131$ one order deeper, while
$9$ are needed deeper still and the remaining four less deeply (two
enter only at order $\eps$, and two multiply coefficients that
vanish identically).  This accounting is settled before any
integration, and it sets the scale of the entire computation.

%======================================================================
\section{Boundary conditions}
\label{sec:boundary}

\subsection{The space of solutions and the conditions upon it}

The path-ordered solution \eqref{eq:pexp} involves the undetermined
vector $\vec F(p_0) = \sum_k \eps^k\,\vec C^{(k)}$: one constant per
master per order.  Almost all of these are fixed by consistency
alone, without evaluating a single new integral.  The most productive
condition --- derived in the next subsection --- is the
\emph{exclusion of local behaviour}: near each edge of the triangle
the differential equation offers local modes with exponents read from
its residue matrices (Section~\ref{sec:local}), the integral
representation permits only some of them, and every offered mode that
the representation forbids annihilates one combination of constants.

Three further conditions complete the analysis.

\emph{Elementary normalizations.}  The unit-power master of the
lowest sector is the phase-space volume itself, a ratio of
$\Gamma$-functions in closed form (Example~\ref{ex:volume}); its
Laurent expansion is the boundary datum of that sector wherever the
sector occurs.

\emph{Inheritance.}  Families are analysed in an order compatible
with the subsector relation.  A master of a subsector, once
evaluated, is literally the same integral in every family in which
that subsector reappears --- the identification of
Section~\ref{sec:classes} is an equality of integrals, not an
analogy.  Expanding the known function at the base point of the new
family and comparing with $\vec I = T\vec F$ there fixes the
corresponding components of $T(p_0)\,\vec F(p_0)$, one inhomogeneous
linear condition on the $\vec C^{(k)}$ per order.

\emph{Regularity at spurious loci.}  A letter may vanish at a locus
interior to the domain in which a given master's integral
representation converges, so that this master is finite there while
the equation still offers modes $u^{\lambda}$ in the distance $u$ to
the locus.  Every offered mode that is singular there --- those with
$\operatorname{Re}\lambda < 0$ for small $\eps$, together with the
powers of $\log u$ accompanying $\lambda = 0$ --- must then have
vanishing coefficient.

\subsection{Power counting at the boundary}
\label{sec:powercounting}

Consider a master given by a convergent cut phase-space integral,
and examine one edge of the triangle, say $w \to 0$.  In the limit
the integral is a sum over momentum regions --- hard, soft,
collinear to one or the other direction --- each of which is a
homogeneous scaling of the integration momenta and therefore
contributes with a definite overall power of $w$.  That power is
counted as follows.  The regulator enters only through the measure,
so the propagators of a region, each scaling as an integer power of
$w$, contribute an integer, while a momentum taken soft,
$k \sim w$, contributes
\begin{equation}
  \dd^d k \;\sim\; w^{\,d} \;=\; w^{\,4-2\eps} ,
\end{equation}
and a momentum taken collinear, $k^+ \sim 1$, $k^- \sim w$,
$k_\perp \sim \sqrt w$, contributes
\begin{equation}
  \dd k^+\,\dd k^-\,\dd^{d-2}k_\perp
  \;\sim\; w^{\,d/2} \;=\; w^{\,2-\eps} .
\end{equation}
The exact function therefore has an expansion in the modes
\begin{equation}
  w^{\,n + a\eps},
  \qquad n \in \Z_{\ge n_0} ,
\end{equation}
in which
\begin{equation}
  a \;=\; -2\,j_{\text{soft}} \;-\; j_{\text{coll}}
  \label{eq:regionexponent}
\end{equation}
counts the momenta the region scales.  The list of admissible $a$ is
therefore short and can be written down without any computation.  The
integrals of the double-real channel have two loop momenta, so
$j_{\text{soft}} + j_{\text{coll}} \le 2$, and the possible pairs
$(j_{\text{soft}}, j_{\text{coll}})$ --- $(0,0)$, $(0,1)$, $(1,0)$
and $(0,2)$, $(1,1)$, $(2,0)$ --- give
\begin{equation}
  a \;\in\; \{\,0,\ -1,\ -2,\ -3,\ -4\,\} ,
\end{equation}
six pairs but five values, since $(1,0)$ and $(0,2)$ both give
$a = -2$; and five is exactly the spectrum of $\eps$-coefficients
observed on the soft edge over the whole channel.  Which of the five
actually occur for a given master is decided by its propagator
content, since that is what decides which regions exist; the datum is
coarse and common to whole groups of sectors.  The differential
equation, on the other hand, offers local modes with exponents drawn
from the eigenvalues of the residue matrix at the corresponding
letter: in the notation of Section~\ref{sec:local}, the general
solution near $w = 0$ with residue matrix $R_0$ organizes itself as
\begin{equation}
  \vec F(v,w) \;=\; \sum_{\lambda}
  w^{\lambda}\, \vec f_{\lambda}(v, \log w) ,
  \qquad
  \lambda \in \eps\operatorname{spec} R_0 + \Z_{\ge 0} ,
  \label{eq:localmodes}
\end{equation}
with each $\vec f_\lambda$ polynomial in $\log w$ and linear in the
boundary constants.

The two lists must be brought into one basis before they are compared,
and they are not yet in one: the region analysis constrains the master
$\vec I = T \vec F$, whereas \eqref{eq:localmodes} lists the exponents
of the canonical $\vec F$, whose integer parts are non-negative by
construction.  If $T$ has order $\nu$ at the edge --- a pole of order
$-\nu$ when $\nu < 0$ --- then the exponents offered for the master are
those of \eqref{eq:localmodes} shifted componentwise by $\nu$,
\begin{equation}
  \lambda_{\vec I} \;\in\;
  \nu + \eps\operatorname{spec}R_0 + \Z_{\ge0} ,
  \label{eq:shiftedmodes}
\end{equation}
and it is this shifted list that the power counting constrains; it is
also where the negative integer parts used in
Section~\ref{sec:endpoint} come from.  Whenever \eqref{eq:shiftedmodes}
offers an exponent that the power counting excludes, the entire
corresponding mode must vanish, $\vec f_{\lambda} \equiv 0$: one exact,
homogeneous, linear condition on the boundary constants per excluded
mode.  The same analysis applies at $v \to 0$ and on the soft edge
$1 - v - w \to 0$, with the distance to the edge as the scaling
variable.

Two cautions are essential in practice.  The admissible exponents must
be derived in the same variables in which the local analysis is
performed, a rationalizing substitution altering the measure by its
Jacobian and shifting the integer parts.  And degenerate (Jordan)
directions of the residue matrix couple modes logarithmically, so the
exclusion argument must be applied to generalized eigenspaces, not to
eigenvalues alone.

\subsection{The kernel of the constraint system}
\label{sec:kernel}

All conditions of the preceding subsections are linear in the
boundary constants of a given block and assemble, order by order in
$\eps$, into one system
\begin{equation}
  C\;\vec C^{(k)} \;=\; \vec d^{\,(k)} ,
\end{equation}
homogeneous ($\vec d = 0$) for the excluded modes and the regularity
conditions, inhomogeneous for the inherited values and the elementary
normalizations.  The structure is the affine-linear one already met at
the completion equation: the solution set is an affine subspace, one
particular solution plus $\ker C$, and the number of genuinely
undetermined constants is $\dim\ker C$ --- a property of the
constraint matrix alone, known before a single integral is
evaluated.  Constants surviving in the kernel are then
identified across blocks: the same subsector appearing in several
families is the same integral, by the equality established in
Section~\ref{sec:classes}, so it contributes one constant and not
several.  For the double-real channel the constraint analysis carried
out this way leaves at most thirty-three candidate periods for the
whole problem, of which the lowest tier is established.

\begin{definition}[Boundary period]
A \emph{boundary period} is the value of a cut phase-space integral
at a distinguished limit of the kinematics, the limits being taken in
a definite order; equivalently, a constant of the general solution
that no linear condition of the kind above determines.
\end{definition}

At the weights relevant here the boundary periods evaluate to
$\Q$-linear combinations of products of $\zeta$-values and kindred
constants.  Each
is established by an exact derivation --- through a Mellin--Barnes
representation, an expansion by momentum regions, or hypergeometric
summation --- with high-precision numerical evaluation used to
corroborate, never to prove.  The Mellin--Barnes representation
rests on one identity, the splitting of a sum in the denominator,
\begin{equation}
  \frac{1}{(X+Y)^{\lambda}}
  \;=\; \frac{1}{\Gamma(\lambda)}\,
  \frac{1}{2\pi i}\int_{-i\infty}^{+i\infty}\!\dd z\;
  \Gamma(\lambda + z)\,\Gamma(-z)\,
  \frac{X^{z}}{Y^{\lambda+z}} ,
  \label{eq:mb}
\end{equation}
the contour separating the poles of $\Gamma(\lambda+z)$ from those of
$\Gamma(-z)$.  Applied to a propagator or to a sum of Feynman
parameters it trades an integration over the integrand for a contour
integral of $\Gamma$-functions, in which the parametric integrals can
be done one at a time.

\begin{example}[A period end to end]
\label{ex:period}
The simplest ordered-limit integral of the whole subject is
\begin{equation}
  P(a,b) \;=\; \int_0^1 \dd x\;
  x^{-1+a\eps}\,(1-x)^{-1+b\eps}
  \;=\; \frac{\Gamma(a\eps)\,\Gamma(b\eps)}
             {\Gamma\big((a+b)\eps\big)} ,
\end{equation}
Euler's beta integral, convergent for $\operatorname{Re}(a\eps)$ and
$\operatorname{Re}(b\eps)$ positive and continued from there.  Its
expansion is obtained from
$\log\Gamma(1+z) = -\gamma_{\mathrm E} z + \sum_{k\ge2}
\frac{(-1)^k\zeta_k}{k}z^k$ after writing
$\Gamma(z) = \Gamma(1+z)/z$:
\begin{equation}
  P(a,b) \;=\; \frac{(a+b)\eps}{a b\,\eps^{2}}\;
  \frac{\Gamma(1+a\eps)\Gamma(1+b\eps)}{\Gamma\big(1+(a+b)\eps\big)}
  \;=\; \frac{a+b}{ab\,\eps}
  \Big[\,1 \;-\; a\,b\,\zeta_2\,\eps^{2} \;+\; O(\eps^{3})\,\Big] ,
\end{equation}
since the terms linear in $\gamma_{\mathrm E}$ cancel between
numerator and denominator and the $\zeta_2$ coefficient is
$\tfrac12\big[a^2+b^2-(a+b)^2\big] = -ab$.  A $\zeta$-value has
appeared at the first nontrivial order, from a $\Gamma$-function and
nothing else; the periods of the calculation are of this character,
one or two integrations deeper.
\end{example}

\begin{example}[The phase-space volume]
\label{ex:volume}
The volume of the massless $n$-particle phase space in $d$
dimensions is fixed by Lorentz invariance and dimensional analysis
up to a constant, and iterating the two-body formula gives it in
closed form:
\begin{equation}
\begin{aligned}
  \int \prod_{i=1}^{n} \dd^d p_i\, \delta^{(+)}(p_i^2)\;
  \delta^d\Big(q - \sum_i p_i\Big)
  \;=\;&
  \big(q^2\big)^{\,(n-1)(d/2-1)-1}
  \left(\frac{\pi^{d/2-1}}{2}\right)^{\!n-1}\\[0.3ex]
  &\times\;
  \frac{\Gamma(d/2-1)^{\,n}}
       {\Gamma\big((n-1)(d/2-1)\big)\,
        \Gamma\big(n(d/2-1)\big)} ,
\end{aligned}
\end{equation}
a pure ratio of $\Gamma$-functions; other normalizations of the
measure differ by powers of $2\pi$ only.  (At $d = 4$ and $n = 3$
this reduces to the familiar $\Phi_3 = \pi^2 q^2/8$.)  It is the
unit-power master of the lowest sector, requires no differential
equation, and supplies through its $\eps$-expansion the exact
$\zeta$-values that inheritance propagates upward.
\end{example}

\begin{example}[The two-dimensional system, concluded]
\label{ex:li2concluded}
The system of Example~\ref{ex:li2} can be solved exactly in $\eps$,
which displays the whole accounting at once.  The first component is
constant, $F_1 = C_1$, and the second obeys
\begin{equation}
  \partial_w F_2 \;=\; -\frac{\eps}{w}\,F_2
  \;-\; \frac{\eps}{1-w}\,C_1 .
\end{equation}
Its homogeneous solution is $w^{-\eps}$, and the particular solution
regular at the origin is found by inserting $F_2 = \sum_{n\ge1}
a_n w^n$ and $-\eps C_1(1-w)^{-1} = -\eps C_1\sum_{m\ge0}w^m$, which
gives $n\,a_n = -\eps\,a_n - \eps\,C_1$, so that
\begin{equation}
  F_2 \;=\; C_2\,w^{-\eps}
  \;-\; \eps\,C_1 \sum_{n\ge1}\frac{w^{\,n}}{n+\eps} .
  \label{eq:li2exact}
\end{equation}
Expanding the sum in $\eps$ reproduces
$\eps\,C_1\log(1-w) + \eps^2 C_1\operatorname{Li}_2(w) + O(\eps^3)$,
in agreement with Example~\ref{ex:li2}.  The two boundary constants
are visible in \eqref{eq:li2exact}: $C_2$ multiplies the mode
$w^{-\eps}$ offered by the eigenvalue $-1$ of the residue matrix at
$w = 0$, and $C_1$ multiplies the mode $w^0$ offered by the
eigenvalue $0$.  The exclusion of Section~\ref{sec:powercounting}
can now be carried out in full.  Suppose the integral represented by
$F_2$ has, in the limit $w \to 0$, only the hard region, so that
\eqref{eq:regionexponent} gives $a = 0$ and its admissible exponents
are the non-negative integers alone.  The equation offers
\begin{equation}
  \eps\operatorname{spec}R_0 + \Z_{\ge0}
  \;=\; \{\,0,\ -\eps\,\} + \Z_{\ge 0} ,
\end{equation}
so the mode $w^{-\eps}$ is offered and forbidden, and its coefficient
must vanish: $C_2 = 0$.  Note that the argument is the comparison of
exponent lists, not an appeal to finiteness --- $w^{-\eps}$ is itself
bounded, and indeed vanishing, as $w \to 0$ for
$\operatorname{Re}\eps < 0$, so finiteness alone would exclude
nothing.  The remaining constant $C_1$ is inherited, the first
component being the volume-type subsector master already known.  The
kernel of the constraint system is trivial, no new period appears,
and $\operatorname{Li}_2(w)$ carries no free constant --- in
miniature, the accounting of the full calculation.
\end{example}

%======================================================================
\section{Behaviour at the edges of phase space}
\label{sec:endpoint}

Section~\ref{sec:local} identified the exponents and the mode
structure at each singular locus, and Section~\ref{sec:boundary}
fixed the constants that multiply them.  One further step converts
the modes into the objects the cross section needs.

Fix an edge of the triangle and let $u$ denote the distance to it:
$u = 1-v-w$ on the soft edge, where this matters most since the Born
configuration sits there, and $u = v$ or $u = w$ on the collinear
edges, where the same identities apply.  Along the soft edge one
variable survives, say $v$, with $w = 1-v-u$.  The master then has,
near that edge, the expansion
\begin{equation}
  \vec I(v,w) \;=\; \sum_{n,a} u^{\,n+a\eps}\;
  \vec g_{\,n,a}(v) \;+\; \dots ,
  \label{eq:edgeexpansion}
\end{equation}
the dots standing for the accompanying powers of $\log u$, and the
cross section integrates it against functions smooth up to the edge.
The coefficients are obtained in two steps.  Projecting the canonical
solution onto the Frobenius basis \eqref{eq:frobenius} gives the
coefficient of each canonical mode: with $R_u$ the residue matrix at
the edge,
\begin{equation}
  \lim_{u\to0}\; \Pi_\lambda\, u^{-\eps R_u}\,\vec F ,
  \label{eq:modeprojection}
\end{equation}
$\Pi_\lambda$ the projector onto the generalized eigenspace of
exponent $\lambda$, the limit existing because $H(u) \to 1$, and for
generic small $\eps$ --- the non-resonance regime of
Section~\ref{sec:local}, $\lvert\operatorname{Re}\eps(\lambda_i -
\lambda_j)\rvert < 1$, so that the modes do not mix under the limit.
The $\vec g_{\,n,a}$ of \eqref{eq:edgeexpansion} are then the
Laurent coefficients of $T$ at the edge contracted with these
canonical mode coefficients, and it is that step which produces the
exponents $n \le -1$: the canonical $\vec F$ has non-negative integer
parts by construction, and the negative ones come from the pole of
$T$, exactly as Section~\ref{sec:powercounting} records.  On the soft
edge the $\vec g_{\,n,a}$ are functions of $v$, not numbers, and it
is those functions which multiply the distributions below.

For such a mode the expansion in $\eps$ and the integration against a
test function do not commute.  The correct object is the
distributional identity
\begin{equation}
  u^{\,-1+a\eps} \;=\;
  \frac{1}{a\eps}\,\delta(u)
  \;+\;
  \sum_{k\ge0} \frac{(a\eps)^k}{k!}
  \left[\frac{\log^k u}{u}\right]_+ ,
  \label{eq:plus}
\end{equation}
where the \emph{plus distribution} is defined by
\begin{equation}
  \int_0^1 \dd u\; \big[g(u)\big]_+\, f(u)
  \;=\; \int_0^1 \dd u\; g(u)\,\big[f(u) - f(0)\big]
\end{equation}
for test functions $f$ smooth at $u = 0$.  The proof of
\eqref{eq:plus} is one line: integrate against $f$, add and
subtract $f(0)$,
\begin{equation}
  \int_0^1 \dd u\;u^{\,-1+a\eps} f(u)
  = \frac{f(0)}{a\eps}
  + \int_0^1 \dd u\;u^{\,-1+a\eps}\big[f(u) - f(0)\big] ,
\end{equation}
where the first term uses $\int_0^1 u^{-1+a\eps}\dd u =
1/(a\eps)$, convergent for $\mathrm{Re}(a\eps) > 0$ and continued
from there; the second integrand may be expanded in $\eps$ term by
term because $f(u) - f(0)$ vanishes at $u = 0$, making every
$u^{-1}\log^k u\,[f(u)-f(0)]$ integrable.  Modes $u^{\,n+a\eps}$ with
$n \le -2$ are treated by the same subtraction carried to higher
order and produce derivatives of $\delta(u)$: for $n = -2$,
subtracting the first \emph{two} Taylor terms of the test function
gives
\begin{equation}
  u^{\,-2+a\eps} \;=\;
  \frac{\delta(u)}{a\eps-1}
  \;-\; \frac{\delta'(u)}{a\eps}
  \;+\; \sum_{k\ge0} \frac{(a\eps)^k}{k!}
  \left[\frac{\log^k u}{u^{2}}\right]_{++} ,
\end{equation}
with the double subtraction defined, in parallel with the single
one, by
\begin{equation}
  \int_0^1 \dd u\; \big[g(u)\big]_{++}\, f(u)
  \;=\; \int_0^1 \dd u\; g(u)\,
  \big[f(u) - f(0) - u\,f'(0)\big] ,
\end{equation}
and $\delta'(u)$ acting as
$\int_0^1 \dd u\,\delta'(u) f(u) = -f'(0)$.  The derivation is the
same one-line computation, now using
$\int_0^1 u^{-2+a\eps}\,\dd u = 1/(a\eps-1)$, convergent for
$\mathrm{Re}(a\eps) > 1$ and continued from there.

As printed, \eqref{eq:plus} is exact on the collinear edges, where $u$
is $v$ or $w$ and the range is $(0,1)$.  On the soft edge the range of
$u = 1-v-w$ at fixed $v$ is $(0,\,1-v)$, and the endpoint enters:
\begin{equation}
  \int_0^{L} \dd u\; u^{\,-1+a\eps}
  \;=\; \frac{L^{\,a\eps}}{a\eps} ,
  \qquad L = 1-v ,
\end{equation}
so that the coefficient of $\delta(u)$ carries a factor
$(1-v)^{a\eps}$, itself to be expanded in the regulator afterwards;
the plus distributions acquire the corresponding rescaling.  The
identity is therefore applied at fixed $v$, with this factor absorbed
into the coefficient functions $\vec g_{\,n,a}(v)$ of
\eqref{eq:edgeexpansion}.

Had one expanded the \emph{mode} first, the left-hand side would have
become the series
\begin{equation}
  \frac{1}{u}\,\sum_{k\ge0} \frac{(a\eps)^k}{k!}\,\log^k u ,
\end{equation}
each term of which is non-integrable at $u = 0$: the
$\delta$-distribution --- the very content the cross section needs
--- is destroyed by premature expansion.  For this reason the
analytic evaluation carries the exact local data
\eqref{eq:edgeexpansion} next to the expanded polylogarithmic form of
each master.  The identity \eqref{eq:plus} is applied to the
unexpanded modes, and only afterwards is everything expanded in the
regulator.  No
resummation is involved anywhere in this step: the residue matrices
supply the exponents, the boundary analysis supplies the
coefficients, and \eqref{eq:plus} converts the two into
distributions.

%======================================================================
\section{Search and proof}
\label{sec:proof}

Nothing in the construction depends on how a transformation was
found.  An ansatz, a numerical fit, arithmetic modulo a prime ---
each produces a candidate and nothing more; what makes the candidate
a theorem is the identity
\begin{equation}
  T^{-1} A\, T \;-\; T^{-1}\dd T
  \;=\; \eps \sum_a R_a\,\dd\log\varphi_a ,
  \label{eq:canonicalid}
\end{equation}
verified exactly in both variables.  The same asymmetry runs through
the whole subject: flatness establishes a connection, and a boundary
period enters only through a derivation in which numerics appear
afterwards, as a check on arithmetic, never as evidence.  Where a
verification is performed modulo primes at random points, its
failure probability is bounded and stated, and the bound is
elementary: a nonzero polynomial of total degree $D$ vanishes at a
point drawn uniformly from $\Fp^{\,n}$ with probability at most
$D/p$ (the Schwartz--Zippel lemma), so a false identity of degree
at most $D$ survives $r$ independent evaluations with probability
at most
\begin{equation}
  \Big(\frac{D}{p}\Big)^{\!r} .
\end{equation}
For a sixty-three--bit prime, $p > 4\cdot 10^{18}$, and identities of
degree $D \le 10^4$, this is $(D/p)^2 < 10^{-29}$ already at
$r = 2$.  It is the stated bound --- not an impression of
reliability --- that the result carries.

It is honest to record the state of the calculation these notes
describe.  At the time of writing, all but the three three-root
families carry canonical forms established in this way ---
eighty-eight of the ninety-one --- while the three-root families
await the Galois-algebraic treatment of
Section~\ref{sec:galois}.  Of the boundary periods, the lowest tier
is established and the higher ones are not; nor is the local mode
analysis of every master at every edge.  The method is settled: each
step above is defined and each of its claims is decided by an exact
identity.  What remains is its execution on the families and the
periods still open.

%======================================================================
\section{Notation and terminology}
\label{sec:summary}

\begin{description}[style=unboxed,leftmargin=0pt,itemsep=0pt,
                    topsep=2pt]
\item[Integral family / representative family.] A complete set of
inverse propagators \eqref{eq:familyint} / its representative modulo
cut-preserving redefinitions of the loop momenta.
\item[Sector, subsector.] The subset of propagators present with
positive power; a subsector arises by deleting lines.
\item[Master integral.] Basis element of a family's integrals modulo
integration-by-parts identities.
\item[Connection, flatness.] The matrix one-form $A$ in
$\dd\vec I = A\,\vec I$; the integrability condition
$\dd A = A \wedge A$.
\item[Block, class.] A strongly connected component of the coupling
graph of a system (a diagonal block of the triangular connection);
its equivalence class under basis permutation composed optionally
with $v \leftrightarrow w$.
\item[Canonical form ($\eps$-form), letters, alphabet, residues.]
The form $A = \eps\sum_a R_a\,\dd\log\varphi_a$ with constant $R_a$;
the functions $\varphi_a$; their set; the matrices $R_a$.
\item[Completion equation.] The linear equation
\eqref{eq:completion} for the entries below the diagonal.
\item[Root content, rationalizing substitution.] The rank of the
group of square classes in an alphabet; a rational change of
variables under which roots become rational.
\item[Galois algebra of the roots.] The extension
$\Q(v,w,\eps)[r_1,\dots,r_k]/(r_i^2 - q_i)$ with automorphism group
$(\Z/2)^k$ acting by sign changes.
\item[Boundary period.] A constant surviving all boundary
conditions; the value of a cut phase-space integral at an ordered
limit.
\item[Local modes.] The exponents $n + a\eps$ and logarithm powers of
the Frobenius solutions at a singular locus, read from the residue
matrix.
\end{description}

\paragraph{Further reading.}  Of the references below, Henn's
paper~\cite{henn} introducing the canonical form is short and repays
being read first, and Duhr's lectures~\cite{duhr} are the standard
entry to iterated integrals and polylogarithms, including the shuffle
and stuffle structure these notes pass over.  The others are the
sources of the individual constructions used above.  Lee's
Libra~\cite{lee} and Meyer's CANONICA~\cite{meyer} are the standard
implementations of the reduction to canonical form.

\addcontentsline{toc}{section}{References}
\sloppy\small
\begin{thebibliography}{99}
\bibitem{kotikov} A.~V. Kotikov, Phys.\ Lett.\ B \textbf{254} (1991)
  158.
\bibitem{gr} T.~Gehrmann and E.~Remiddi, Nucl.\ Phys.\ B \textbf{580}
  (2000) 485; hep-ph/9912329.
\bibitem{henn} J.~M. Henn, Phys.\ Rev.\ Lett.\ \textbf{110} (2013)
  251601; arXiv:1304.1806.
\bibitem{am} C.~Anastasiou and K.~Melnikov, Nucl.\ Phys.\ B
  \textbf{646} (2002) 220; hep-ph/0207004.
\bibitem{duhr} C.~Duhr, arXiv:1411.7538.
\bibitem{moser} J.~Moser, Math.\ Z.\ \textbf{72} (1959) 379.
\bibitem{lee} R.~N. Lee, JHEP \textbf{04} (2015) 108;
  arXiv:1411.0911.
  R.~N. Lee, Comput.\ Phys.\ Commun.\ \textbf{267} (2021) 108058;
  arXiv:2012.00279 (Libra).
\bibitem{meyer} C.~Meyer, JHEP \textbf{04} (2017) 006;
  arXiv:1611.01087.
  C.~Meyer, Comput.\ Phys.\ Commun.\ \textbf{222} (2018) 295;
  arXiv:1705.06252 (CANONICA).
\bibitem{pak} A.~Pak, J.\ Phys.\ Conf.\ Ser.\ \textbf{368} (2012)
  012049; arXiv:1111.0868.
\bibitem{besier} M.~Besier, D.~van Straten and S.~Weinzierl,
  arXiv:1809.10983.
  M.~Besier, P.~Wasser and S.~Weinzierl, Comput.\ Phys.\ Commun.\
  \textbf{253} (2020) 107197; arXiv:1910.13251.
\bibitem{chen} K.-T. Chen, Bull.\ Amer.\ Math.\ Soc.\ \textbf{83}
  (1977) 831.
\bibitem{goncharov} A.~B. Goncharov, arXiv:math/0103059.
\bibitem{beauville} A.~Beauville, Complex Algebraic Surfaces, 2nd
  ed., London Mathematical Society Student Texts 34, Cambridge
  University Press (1996).
\end{thebibliography}

\end{document}
